\documentclass[11pt,twoside]{book}

%%Page Size (rev. 08/19/2016)
%\usepackage[inner=0.5in, outer=0.5in, top=0.5in, bottom=0.5in, papersize={6in,9in}, head=12pt, headheight=30pt, headsep=5pt]{geometry}
\usepackage[inner=0.6in, outer=0.4in, top=0.5in, bottom=0.5in, papersize={5.5in,8.5in}, head=12pt, headheight=30pt, headsep=5pt]{geometry}
%% width of textblock = 324 pt / 4.5in
%% A5 = 5.8 x 8.3 inches -- if papersize is A5, then margins should be [inner=0.75in, outer=0.55in, top=0.4in, bottom=0.4in]


%%Header (rev. 4/11/2011)
\usepackage{fancyhdr}
 \pagestyle{fancy}
\renewcommand{\chaptermark}[1]{\markboth{#1}{}}
\renewcommand{\sectionmark}[1]{\markright{#1}}
 \fancyhf{}
\fancyhead[LE,RO]{\thepage}
\fancyhead[CE]{\leftmark}
\fancyhead[CO]{\rightmark}
 \fancypagestyle{plain}{ %
\fancyhf{} % remove everything
\renewcommand{\headrulewidth}{0pt} % remove lines as well
\renewcommand{\footrulewidth}{0pt}}

\fancypagestyle{centernumber}{ %
\fancyhf{} % remove everything
\cfoot{\thepage}
\renewcommand{\headrulewidth}{0pt} % remove lines as well
\renewcommand{\footrulewidth}{0pt}
}

\fancypagestyle{plain}{%
\fancyhf{}    %pagestyle definition
\renewcommand{\headrulewidth}{0pt} % remove lines as well
\renewcommand{\footrulewidth}{0pt}
}

%%Blank page before new chapter
\makeatletter
\def\cleardoublepage{\clearpage\if@twoside \ifodd\c@page\else
\hbox{}
\vspace*{\fill}
\begin{center}
\phantom{I gave you all the clues.}
\end{center}
\vspace{\fill}
\thispagestyle{empty}
\newpage
\if@twocolumn\hbox{}\newpage\fi\fi\fi}
\makeatother


\usepackage[autocompile,allowdeprecated=false]{gregoriotex}
\usepackage{gregoriosyms}
\gresetgregoriofont[op]{greciliae}

%%Titles (rev. 9/4/2011) -- TOCLESS --- lets you have sections that don't appear in the table of contents

\setcounter{tocdepth}{0}
\setcounter{secnumdepth}{-1}

%%Titlesec modified for Triduum booklet (3.22.20)
\usepackage[compact,nobottomtitles*]{titlesec}
\titlespacing*{\chapter}{0pt}{-30pt}{0pt}
\titlespacing*{\section}{0pt}{*0}{*1}
\titlespacing*{\subsubsection}{0pt}{*1}{*1}
\titlespacing*{\subsubsubsection}{0pt}{10pt}{*0}
\titleformat{\part} {\normalfont\Huge\sc\center}{\thechapter}{1em}{}
\titleformat{\chapter} {\normalfont\LARGE\sc\center}{\thechapter}{1em}{}
\titleformat{\section} {\normalfont\Large\sc\center}{\thesection}{1em}{}
\titleformat{\subsection} {\normalfont\Large\sc\center}{\thesubsubsection}{1em}{}
\titleformat{\subsubsection}{\normalfont\large\sc\center}{\thesubsubsubsection}{1em}{}
\titleformat{\paragraph}{\normalfont\normalsize\sc\center}{\thesubsubsubsection}{1em}{}

\newcommand{\nocontentsline}[3]{}
\newcommand{\tocless}[2]{\bgroup\let\addcontentsline=\nocontentsline#1{#2}\egroup} %% lets you have sections that don't appear in the table of contents


%%%

%%Index (rev. December 11, 2013)
\usepackage[noautomatic,nonewpage]{imakeidx}


\makeindex[name=incipit,title=Index]
\indexsetup{level=\section,toclevel=section,noclearpage}

\usepackage[indentunit=8pt,rule=.5pt,columns=2]{idxlayout}


%%Table of Contents (rev. May 16, 2011)

%\usepackage{multicol}
%\usepackage{ifthen}
%\usepackage[toc]{multitoc}

%% General settings (rev. January 19, 2015)

\usepackage[normalem]{ulem}

\usepackage[latin,german]{babel}
\usepackage{lettrine}

\usepackage{paracol}

\usepackage{fontspec}

\setmainfont[Ligatures=TeX,BoldFont=MinionPro-Bold,ItalicFont=MinionPro-It, BoldItalicFont=MinionPro-BoldIt]{MinionPro-Regular-Modified.otf}

%% Style for translation line
\grechangestyle{translation}{\fontsize{10}{10}\it\selectfont}
\grechangestyle{annotation}{\fontsize{10}{10}\selectfont}
\grechangestyle{commentary}{\textnormal\selectfont}
\gresetcustosalteration{invisible}

%\grechangedim{annotationseparation}{0.1cm}{scalable}

%\GreLoadSpaceConf{smith-four}

\frenchspacing

\usepackage{indentfirst} %%%indents first line after a section

\usepackage{graphicx}
%\usepackage{tocloft}

%%Hyperref (rev. August 20, 2011)
%\usepackage[colorlinks=false,hyperindex=true,bookmarks=true]{hyperref}
\usepackage{hyperref}
\hypersetup{pdftitle={Vesperale O.P. 2017}}
\hypersetup{pdfauthor={Order of Preachers}}
\hypersetup{pdfsubject={Liturgy}}
\hypersetup{pdfkeywords={Dominican, Liturgy, Order of Preachers, Dominican Rite, Liturgia Horarum, Divine Office}}

\newlength{\drop}



\begin{document}


\raggedbottom

\newcommand{\lectio}[3]{%
  \makebox[0pt][l]{#1}%
  \makebox[\textwidth][c]{#2}%
  \makebox[0pt][r]{\normalsize{\textnormal{#3}}}}

  %%InitialPage
  \begin{titlepage}
      \drop=0.1\textheight
      \centering
      \vspace*{\baselineskip}
     % \rule{\textwidth}{1.6pt}\vspace*{-\baselineskip}\vspace*{2pt}
      %\rule{\textwidth}{0.4pt}\\[\baselineskip]
      {\fontsize{45pt}{45}\selectfont \textsc{Triduum\\ante Pascha}
     % \rule{\textwidth}{0.4pt}\vspace*{-\baselineskip}\vspace{3.2pt}
      %\rule{\textwidth}{1.6pt}\\[\baselineskip]

      \vspace*{100pt}

      \begin{figure}[h]
  \centering
        \includegraphics[width=1.75in]{op-shield.jpg} %
      \end{figure}


      \vfill
      {\fontsize{13pt}{20pt}\selectfont \textsc{München} \par \textsc{Dominikanerkonvent St. Kajetan} \par \scshape{2020} \par }}
    \end{titlepage}




  %%Combination
  \newgeometry{inner=1.5in, outer=1.5in, top=2in, bottom=0.5in}

    \tableofcontents

  \restoregeometry

  \chapter{Trauermette am Gründonnerstag}
  \section{Officium Lectionis}
      \index[Varia]{Herr offne} \label{Herr offne (Varia)} \grecommentary[0pt]{} \gresetinitiallines{1} \grechangestyle{initial}{\fontsize{36}{36}\selectfont} \grechangedim{maxbaroffsettextleft@nobar}{12 cm}{scalable} \grechangedim{spaceabovelines}{0.5cm}{scalable} \gresetlyriccentering{syllable}  \grechangedim{maxbaroffsettextleft}{0 cm}{scalable} \gregorioscore{chants/herr_offne}
   \subsubsection{Invitatorium}   \index[Invitatorium]{Invitatorium} \label{Invitatorium (Invitatorium)}         \vspace{5pt} \par \gresetlyriccentering{syllable}

\emph{Cantor:}~
\nopagebreak
\greannotation{I}
\gabcsnippet{(c4) Im(d) Kreuz(h_) Je(g)su(h) Chri(j_)sti(h_) fin(g)den(f) wir(e_) Heil.(d_) (::)} \vspace{5pt}

\emph{Alle:}~
\nopagebreak
\gabcsnippet{(c4) Im(d) Kreuz(h_) Je(g)su(h) Chri(j_)sti(h_) fin(g)den(f) wir(e_) Heil.(d_) (::)} \vspace{5pt}

\emph{Cantor:}~
\nopagebreak
\gabcsnippet{Kommt(dh_), lasst(h) uns(h) ju(i)beln(h) vor(h) dem(g) Herrn(h) (,) und(j) zu(h)jauch(h)zen(g) dem(f) Fels(g) uns(g)res(h) Hei(f)les! (f) (:) 
Lasst(f) uns(f) mit(g) Lob(ixi_) sei(h)-nem(h) An(h)ge(h)sicht(f) na(gh)hen, (g) (,) vor(e) ihm(e) jau(g)chzen(g) mit(h) Lie(fe)dern! (d) <sp>R/</sp> (::)} \vspace{5pt}


\emph{Cantor:}~
\nopagebreak
\gabcsnippet{(c4)Denn(d) der(d) Herr(h_) ist(h) ein(h) gro(i)ßer(g) Gott,(h) (,) ein(h) gro(h)ßer(j) Kö(hi)nig(h) üb(h)er(h) all(g)en(h) Gö(f)ttern.(f) (:)  In(d) sei(f)ner(f) Hand(g_) sind(g) die(g) Tie(h)fen(g) der(h) Er(f)de,(g) (,) sein(f) sind(f) die(f) Gi(g)pfel(g) der(f) Ber(gh)ge.(h) (:) Sein(h) ist(g) das(h) Meer,(ixi_) das(h) er(h) ge(f)macht(g) hat,(g) (,) das(g) tro(h)cke(g)ne(f) Land,(g_) das(g) sei(g)ne(e) Hän(g)de(g)  ge(h)bil(fe)det.(d) (::)} \vspace{5pt}

\emph{Alle:}~
\nopagebreak
\gabcsnippet{(c4) Im(d) Kreuz(h_) Je(g)su(h) Chri(j_)sti(h_) fin(g)den(f) wir(e_) Heil.(d_) (::)} \vspace{5pt}

\emph{Cantor:}~
\nopagebreak

\gabcsnippet{(c4)Kommt(d_), lasst(d) uns(d) nie(h)der(h)fal(h)len(h), uns(h) vor(h) ihm(i) ver(g)nei(hi)gen,(h) (,) lasst(h) uns(j) nie(h)der(h)knien(h) vor(g) dem(f) Herrn,(g) un(g)serm(h) Schö(f)pfer!(f) (:) Denn(f) er(g) ist(g) un(f)ser(g) Gott,(h) (,) wir(h) sind(h) das(h) Volk(ixi_) sei(h)ner(f) Wei(gh)de,(g) (,) die(g) Her(fg)de,(g) von(e) sei(g)ner(h) Hand(f) ge(e)führt.(d) 
(::)}

 \vspace{5pt}


\emph{Alle:}~
\nopagebreak
\gabcsnippet{(c4) Im(d) Kreuz(h_) Je(g)su(h) Chri(j_)sti(h_) fin(g)den(f) wir(e_) Heil.(d_) (::)}

 \vspace{5pt}

\emph{Cantor:}~
\nopagebreak
\gabcsnippet{(c4) Ach,(d_) wür(d)det(d) ihr(d) doch(d) heu(dh)te(h) auf(h) sei(h)ne(h) Sti(i)mme(g) hö(hi)ren!(h) (;) „Ver(h)här(h)tet(h) eu(h)er(h) Herz(g) nicht(h) wie(h) in(h) Me(f)rí(g)ba,(g) (,) wie(h) in(h) der(j) Wü(hi)ste(h) am(h) Tag(g) von(h) Ma(f)ssa!(f) (:) Dort(f) ha(f)ben(f) eu(f)re(f) Vä(g)ter(g) mich(f) ver(g)sucht,(h) (,) sie(h) ha(h)ben(h) mich(h) auf(h) die(h) Pro(ixi)be(h) ge(f)stellt(g) und(g) hat(g)ten(g) doch(g) mein(e) Tun(g) ge(h)se(fe)hen.(d) 
(::)}

 \vspace{5pt}


\emph{Alle:}~
\nopagebreak
\gabcsnippet{(c4) Im(d) Kreuz(h_) Je(g)su(h) Chri(j_)sti(h_) fin(g)den(f) wir(e_) Heil.(d_) (::)}

 \vspace{5pt}

\emph{Cantor:}~\nopagebreak

\gabcsnippet{(c4) Vier(d)zig(d) Jah(dh)re(h) war(h) mir(h) dies(h) Ge(h)schlecht(i) zu(g)wi(hi)der,(h) (;) und(h) ich(f) sa(gh)gte:(g) Sie(h) sind(h) ein(j) Volk(h), de(h)ssen(h) Herz(h) in(g) die(f) Ir(g)re(g) geht;(g) (,) denn(h) mei(h)ne(j) We(hi)ge(h) ken(h)nen(g) sie(h) nicht.(f) (:) Da(f)rum(f) ha(f)be(f) ich(f) in(f) mei(f)nem(f) Zorn(g) ge(f)schwo(gh)ren:(h) (,) Sie(h) sol(ixi)len(h) nicht(f) kom(gh)men(g) in(g) das(e) Land(g) mei(g)ner(h) Ru(fe)he.”(d) (::)}

 \vspace{5pt}


\emph{Alle:}~\nopagebreak

\gabcsnippet{(c4) Im(d) Kreuz(h_) Je(g)su(h) Chri(j_)sti(h_) fin(g)den(f) wir(e_) Heil.(d_) (::)}

 \vspace{5pt}

\emph{Cantor:}~
\nopagebreak

\gabcsnippet{(c4) Eh(d)re(d) dem(d) Va(dh)ter,(h) (,) Eh(i)re(h) dem(g) Sohn,(h) (,) Eh(h)re(h) dem(j) Hei(h)li(g)gen(h) Geist.(f) (:) Wie(f) im(g) An(ixi)fang,(h) (,) so(h) auch(h) jetzt(h) und(f) al(g)le(h) Zeit(g) und(g) in(e) E(g)wig(g)keit.(h) A(fvED)men.(d) (::)}

 \vspace{5pt}


\emph{Alle:}~
\nopagebreak
\gabcsnippet{(c4) Im(d) Kreuz(h_) Je(g)su(h) Chri(j_)sti(h_) fin(g)den(f) wir(e_) Heil.(d_) (::)}
  \subsubsection{Hymnus}  \greannotation{I} \index[Hymnus]{Hymnus} \label{Hymnus (Hymnus)} \grecommentary[0pt]{} \gresetinitiallines{1} \grechangestyle{initial}{\fontsize{36}{36}\selectfont} \grechangedim{maxbaroffsettextleft@nobar}{12 cm}{scalable} \grechangedim{spaceabovelines}{0.5cm}{scalable} \gresetlyriccentering{syllable}   

\gabcsnippet{
(c4) Heil(d)ig(e) Kreuz,(gh) du(hg) Baum(h) der(j) Treu(ji)e,(g) (,)
ed(i)ler(j) Baum,(kl) dem(hg) kein(j)er(i) gleich,(h) (;)
kein(hj)er(i) so(g) an(e) Laub(f) und(dc) Blü(de)te,(e) (,)
kein(d)er(h) so(hg) an(ed) Früch(f)ten(e) reich:(d) (;)
Sü(hj)ßes(i) Holz,(g) o(e) sü(f)ße(dc) Nä(de)gel,(e) (,)
wel(d)che(h) sü(hg)ße(ed) Last(f) an(e) euch.(d) (::)}

\gresetinitiallines{0}

\vspace{7pt}

\gabcsnippet{2. Beu(d)ge,(e) ho(gh)her(hg) Baum,(h) die(j) Zwei(ji)ge,(g) (,)
wer(i)de(j) weich(kl) an(hg) Stamm(j) und(i) Ast,(h) (;)
denn(hj) dein(i) har(g)tes(e) Holz(f) muß(dc) tra(de)gen(e) (,)
ei(d)ne(h) kö(hg)nig(ed)li(f)che(e) Last,(d) (;)
gib(hj) den(i) Glie(g)dern(e) dei(f)nes(dc) Schöp(de)fers(e) (,)
an(d) dem(h) Stam(hg)me(ed) lin(f)de(e) Rast.(d) (::)}

\vspace{7pt}

\gabcsnippet{3. Du(d) al(e)lein(gh) warst(hg) wert,(h) zu(j) tra(ji)gen(g) (,)
al(i)ler(j) Sün(kl)den(hg) Lö(j)se(i)geld,(h) (;)
du,(hj) die(i) Plan(g)ke,(e) die(f) uns(dc) ret(de)tet(e) (,)
aus(d) dem(h) Schiff(hg)bruch(ed) die(f)ser(e) Welt.(d) (;)
Du,(hj) ge(i)salbt(g) vom(e) Blut(f) des(dc) Lam(de)mes,(e) (,)
Pfos(d)ten,(h) der(hg) den(ed) Tod(f) ab(e)hält.(d) (::)}

\vspace{7pt}

\gabcsnippet{4. Lob(d) und(e) Ruhm(gh) sei(hg) oh(h)ne(j) En(ji)de(g) (,)
Gott,(i) dem(j) höchs(kl)ten(hg) Herrn,(j) ge(i)weiht.(h) (;)
Preis(hj) dem(i) Va(g)ter(e) und(f) dem(dc) Soh(de)ne(e) (,)
und(d) dem(h) Geist(hg) der(ed) Hei(f)lig(e)keit.(d) (;)
Ei(hj)nen(i) Gott(g) in(e) drei(f) Per(dc)so(de)nen(e) (,)
lo(d)be(h) al(hg)le(ed) Welt(f) und(e) Zeit.(d) (::)
A(ded)men.(cd) (::)}                 \newpage
  \subsection{Psalmodie} \noindent \textsc{1 Ant.} Du hast uns gerettet, Herr, wir preisen deinen Namen auf ewig.

\subsubsection{Psalm 44 (43)} \paragraph{I}

\noindent Gott, wir hörten es mit eigenen Ohren,~$\star$~\nopagebreak

unsere Väter erzählten uns

\noindent von dem Werk, das du in ihren Tagen vollbracht hast,~$\star$~\nopagebreak

in den Tagen der Vorzeit.

\noindent Mit eigener Hand hast du Völker vertrieben,~$\star$~\nopagebreak

sie aber eingepflanzt.

\noindent Du hast Nationen zerschlagen,~$\star$~\nopagebreak

sie aber ausgesät.

\noindent Denn sie gewannen das Land nicht mit ihrem Schwert,~$\star$~\nopagebreak

noch verschaffte ihr Arm ihnen den Sieg;

\noindent nein, deine Rechte war es, dein Arm und dein leuchtendes Angesicht;~$\star$~\nopagebreak

denn du hattest an ihnen Gefallen.

\noindent Du, mein König und mein Gott,~$\star$~\nopagebreak

du bist es, der Jakob den Sieg verleiht.

\noindent Mit dir stoßen wir unsere Bedränger nieder,~$\star$~\nopagebreak

in deinem Namen zertreten wir unsere Gegner.

\noindent Denn ich verlasse mich nicht auf meinen Bogen,~$\star$~\nopagebreak

noch kann mein Schwert mir helfen;

\noindent nein, du hast uns vor unsern Bedrängern gerettet;~$\star$~\nopagebreak

alle, die uns hassen, bedeckst du mit Schande.

\noindent Wir rühmen uns Gottes den ganzen Tag~$\star$~\nopagebreak

und preisen deinen Namen auf ewig.

\noindent Ehre sei dem Vater und dem Sohn~$\star$~\nopagebreak

und dem Heiligen Geist.

\noindent Wie im Anfang, so auch jetzt und alle Zeit~$\star$~\nopagebreak

und in Ewigkeit. Amen.

\vspace{10pt}

\noindent \textsc{Ant.} Du hast uns gerettet, Herr, wir preisen deinen Namen auf ewig. 

%\vspace{10pt}

\newpage

\noindent \textsc{2 Ant.} Verschone dein Volk, o Herr; gib dein Erbe nicht der Schande preis.

\paragraph{II}

\noindent Doch nun hast du uns verstoßen und mit Schmach bedeckt,~$\star$~\nopagebreak

du ziehst nicht mit unserm Heer in den Kampf. 

 \noindent Du lässt uns vor unsern Bedrängern fliehen~$\star$~\nopagebreak

und Menschen, die uns hassen, plündern uns aus. 

 \noindent Du gibst uns preis wie Schlachtvieh,~$\star$~\nopagebreak

unter die Völker zerstreust du uns. 

 \noindent Du verkaufst dein Volk um ein Spottgeld~$\star$~\nopagebreak

und hast an dem Erlös keinen Gewinn. 

 \noindent Du machst uns zum Schimpf für die Nachbarn,~$\star$~\nopagebreak

zu Spott und Hohn bei allen, die rings um uns wohnen.

\noindent Du machst uns zum Spottlied der Völker,~$\star$~\nopagebreak

die Heiden zeigen uns nichts als Verachtung.

\noindent Meine Schmach steht mir allzeit vor Augen~$\star$~\nopagebreak

und Scham bedeckt mein Gesicht

\noindent wegen der Worte des lästernden Spötters,~$\star$~\nopagebreak

wegen der rachgierigen Blicke des Feindes.

\noindent Ehre sei dem Vater und dem Sohn~$\star$~\nopagebreak

und dem Heiligen Geist.

\noindent Wie im Anfang, so auch jetzt und alle Zeit~$\star$~\nopagebreak

und in Ewigkeit. Amen.

\vspace{10pt}

\noindent \textsc{Ant.} Verschone dein Volk, o Herr; gib dein Erbe nicht der Schande preis. \vspace{10pt}

\noindent \textsc{3 Ant.} Steh auf und hilf uns, Herr; in deiner Huld erlöse uns.

\paragraph{III}

\noindent Das alles ist über uns gekommen~†~\nopagebreak

und doch haben wir dich nicht vergessen,~$\star$~\nopagebreak

uns von deinem Bund nicht treulos abgewandt.

\noindent Unser Herz ist nicht von dir gewichen,~$\star$~\nopagebreak

noch hat unser Schritt deinen Pfad verlassen.

\noindent Doch du hast uns verstoßen an den Ort der Schakale~$\star$~\nopagebreak

und uns bedeckt mit Finsternis.

\noindent Hätten wir den Namen unseres Gottes vergessen~$\star$~\nopagebreak

und zu einem fremden Gott die Hände erhoben,

\noindent würde Gott das nicht ergründen?~$\star$~\nopagebreak

Denn er kennt die heimlichen Gedanken des Herzens.

\noindent Nein, um deinetwillen werden wir getötet Tag für Tag,~$\star$~\nopagebreak

behandelt wie Schafe, die man zum Schlachten bestimmt hat. 

 \noindent Wach auf! Warum schläfst du, Herr?~$\star$~\nopagebreak

Erwache, verstoß nicht für immer! 

 \noindent Warum verbirgst du dein Gesicht,~$\star$~\nopagebreak

vergisst unsere Not und Bedrängnis? 

 \noindent Unsere Seele ist in den Staub hinabgebeugt,~$\star$~\nopagebreak

unser Leib liegt am Boden.

\noindent Steh auf und hilf uns!~$\star$~\nopagebreak

In deiner Huld erlöse uns!

\noindent Ehre sei dem Vater und dem Sohn~$\star$~\nopagebreak

und dem Heiligen Geist.

\noindent Wie im Anfang, so auch jetzt und alle Zeit~$\star$~\nopagebreak

und in Ewigkeit. Amen.

\vspace{10pt}

\noindent \textsc{Ant.} Steh auf und hilf uns, Herr; in deiner Huld erlöse uns.

   \subsubsection{Versiculum}  \greannotation{} \index[Versiculum]{Versiculum} \label{Versiculum (Versiculum)} \grecommentary[0pt]{} \gresetinitiallines{0} \grechangestyle{initial}{\fontsize{36}{36}\selectfont} \grechangedim{maxbaroffsettextleft@nobar}{12 cm}{scalable} \grechangedim{spaceabovelines}{0.5cm}{scalable} \gresetlyriccentering{syllable}   \gregorioscore{chants/versiculum_thursday}    \newpage
  \subsection{Lesungen}
   \subsubsection{Erste Lesung} \vspace{5pt} \emph{Anfang der Klagelieder des Propheten Jeremia. Aleph. Weh, wie einsam sitzt da die einst so volkreiche Stadt. Einer Witwe wurde gleich die Große unter den Völkern. Die Fürstin über die Länder ist zur Fron erniedrigt. Beth. Sie weint und weint des Nachts, Tränen auf ihren Wangen. Keinen hat sie als Tröster von all ihren Geliebten. Untreu sind all ihre Freunde, sie sind ihr zu Feinden geworden. Ghimel. Gefangen ist Juda im Elend, in harter Knechtschaft. Nun weilt sie unter den Völkern und findet nicht Ruhe. All ihre Verfolger holten sie ein mitten in der Bedrängnis. - Jerusalem, Jerusalem, bekehre dich zum Herrn, Deinem Gott.} \vspace{5pt} \greannotation{} \index[Erste Lesung]{Erste Lesung} \label{Erste Lesung (Erste Lesung)} \grecommentary[0pt]{Klgl 1, 1--3} \gresetinitiallines{1} \grechangestyle{initial}{\fontsize{36}{36}\selectfont} \grechangedim{maxbaroffsettextleft@nobar}{12 cm}{scalable} \grechangedim{spaceabovelines}{0.5cm}{scalable} \gresetlyriccentering{vowel}   \gregorioscore{chants/va--incipit_lamentatio_ieremiae_prophetae--dominican}
   \subsubsection{II} \vspace{5pt} \emph{Daleth. Die Wege nach Zion trauern, niemand pilgert zum Fest, verödet sind all ihre Tore. Ihre Priester seufzen, ihre Jungfrauen sind voll Gram, sie selbst trägt Weh und Kummer. He. Ihre Bedränger sind an der Macht, ihre Feinde im Glück. Denn Trübsal hat der Herr ihr gesandt wegen ihrer vielen Sünden. Ihre Kinder zogen fort, gefangen, vor dem Bedränger. Gewichen ist von der Tochter Zion all ihre Pracht. Vau. Ihre Fürsten sind wie Hirsche geworden, die keine Weide finden. Kraftlos zogen sie dahin vor ihren Verfolgern. - Jerusalem, Jerusalem, bekehre dich zum Herrn, deinem Gott.} \vspace{5pt} \greannotation{} \index[II]{II} \label{II (II)} \grecommentary[0pt]{Klgl 1, 4--6} \gresetinitiallines{1} \grechangestyle{initial}{\fontsize{36}{36}\selectfont} \grechangedim{maxbaroffsettextleft@nobar}{12 cm}{scalable} \grechangedim{spaceabovelines}{0.5cm}{scalable} \gresetlyriccentering{vowel}   \gregorioscore{chants/va--daleth_viae_sion_lugent--dominican}
   \subsubsection{III} \vspace{5pt} \emph{Zain. Jerusalem denkt an die Tage ihres Elends, ihrer Unrast, an all ihre Kostbarkeiten, die sie einst besessen, als ihr Volk in Feindeshand fiel und keiner ihr beistand. Die Feinde sahen sie an, lachten über ihre Vernichtung. Heth. Schwer gesündigt hatte Jerusalem, deshalb ist sie zum Abscheu geworden. All ihre Verehrer verachten sie, weil sie ihre Blöße gesehen. Sie selbst aber seufzt und wendet sich ab (von ihnen). Teth. Ihre Unreinheit klebt an ihrer Schleppe, ihr Ende bedachte sie nicht. Entsetzlich ist sie gesunken, keinen hat sie als Tröster. Sieh doch mein Elend, o Herr, denn die Feinde prahlen. - Jerusalem, Jerusalem, bekehre dich zum Herrn, deinem Gott.}  \vspace{5pt} \greannotation{} \index[III]{III} \label{III (III)} \grecommentary[0pt]{Klgl 1, 7--9} \gresetinitiallines{1} \grechangestyle{initial}{\fontsize{36}{36}\selectfont} \grechangedim{maxbaroffsettextleft@nobar}{12 cm}{scalable} \grechangedim{spaceabovelines}{0.5cm}{scalable} \gresetlyriccentering{vowel}   \gregorioscore{chants/va--zain_recordata_est_ierusalem--dominican}    \newpage
   \subsubsection{Zweite Lesung}             \vspace{5pt}

Meliton von Sardes († vor 190), Aus einer Osterpredigt

\vspace{5pt}

%\par \hfill(Nn. 65071: SC 123, 95-101)

\lettrine[lines=3]{D}{}ie Propheten haben vieles vorausverkündigt über das Paschamysterium, das Christus ist, „dem Ehre sei in alle Ewigkeit. Amen“. Er kam vom Himmel auf die Erde wegen des leidenden Menschen; den leidenden Menschen zog er wie ein Kleid an im Schoß der Jungfrau und ging hervor als Mensch; durch einen Leib, der dem Leiden ausgesetzt war, nahm er die Leiden des leidenden Menschen auf sich und vernichtete die Leiden des Fleisches. Durch den Geist aber, der nicht sterben konnte, tötete er den Mörder Tod.

Er wurde zum Schlachten geführt wie ein Lamm und getötet wie ein Schaf. Wie aus einem Ägypten erlöste er uns aus dem Dienst der Welt. Er rettete uns aus der Knechtschaft des Teufels wie aus der Hand des Pharaos; er besiegelte unsere Seelen mit seinem eigenen Geist und die Glieder unseres Leibes mit seinem Blut. Er ist es, der Verwirrung über den Tod brachte und den Teufel in Trauer versetzte wie Mose den Pharao. Er schlug die Bosheit und verdammte die Ungerechtigkeit zur Unfruchtbarkeit wie Mose Ägypten.

Er ist es, der uns der Knechtschaft entrissen und uns befreit hat, der uns aus der Finsternis zum Licht führte, vom Tod zum Leben, von der Gewaltherrschaft zu ewigem Königtum, der uns zu einer neuen Priesterschaft machte, zu einem erwählten und ewigen Volk. Er ist das Paschalamm unseres Heils. Er ertrug in vielen vieles: Er wurde in Abel gemordet. 

In Isaak wurden ihm die Füße gefesselt, in Jakob mußte er auswandern. In Josef wurde er verkauft, in Mose ausgesetzt, im Lamm geschlachtet, in David verfolgt, in den Propheten geschmäht. Er wurde Mensch in der Jungfrau, ans Holz gehängt, in das Grab der Erde gesenkt. Er erstand von den Toten und stieg empor zur Höhe des Himmels. Er ist das Lamm, das verstummt, aus der Herde geholt, zum Schlachten geführt, am Abend geopfert und in der Nacht begraben. Am Holz zerbrach man ihn nicht und im Grab verweste er nicht. Er stand von den Toten auf und erweckte den Menschen aus dem Grab der Unterwelt.    \newpage
  \section{Laudes}
  \subsubsection{Psalmodie} \noindent \textsc{1 Ant.} Sieh her, mein Gott, verbirg nicht dein Gesicht, denn mir ist angst, erhöre mich bald.

\subsubsection{Psalm 80}

\noindent Du Hirte Israels, höre,~$\star$~\nopagebreak

der du Josef weidest wie eine Herde!

\noindent Der du auf den Kerubim thronst, erscheine~$\star$~\nopagebreak

vor Efraim, Benjamin und Manasse!

\noindent Biete deine gewaltige Macht auf~$\star$~\nopagebreak

und komm uns zu Hilfe!

\noindent Gott, richte uns wieder auf!~$\star$~\nopagebreak

Lass dein Angesicht leuchten, dann ist uns geholfen.

\noindent Herr, Gott der Heerscharen, wie lange noch zürnst du,~$\star$~\nopagebreak

während dein Volk zu dir betet?

\noindent Du hast sie gespeist mit Tränenbrot,~$\star$~\nopagebreak

sie überreich getränkt mit Tränen.

\noindent Du machst uns zum Spielball der Nachbarn~$\star$~\nopagebreak

und unsere Feinde verspotten uns.

\noindent Gott der Heerscharen, richte uns wieder auf!~$\star$~\nopagebreak

Lass dein Angesicht leuchten, dann ist uns geholfen.

\noindent Du hobst in Ägypten einen Weinstock aus,~$\star$~\nopagebreak

du hast Völker vertrieben, ihn aber eingepflanzt.

\noindent Du schufst ihm weiten Raum;~$\star$~\nopagebreak

er hat Wurzeln geschlagen und das ganze Land erfüllt.

\noindent Sein Schatten bedeckte die Berge,~$\star$~\nopagebreak

seine Zweige die Zedern Gottes.

\noindent Seine Ranken trieb er hin bis zum Meer~$\star$~\nopagebreak

und seine Schößlinge bis zum Eufrat.

\noindent Warum rissest du seine Mauern ein?~$\star$~\nopagebreak

Alle, die des Weges kommen, plündern ihn aus.

\noindent Der Eber aus dem Wald wühlt ihn um,~$\star$~\nopagebreak

die Tiere des Feldes fressen ihn ab.

\noindent Gott der Heerscharen, wende dich uns wieder zu!~$\star$~\nopagebreak

Blick vom Himmel herab, und sieh auf uns!

\noindent Sorge für diesen Weinstock~$\star$~\nopagebreak

und für den Garten, den deine Rechte gepflanzt hat.

\noindent Die ihn im Feuer verbrannten wie Kehricht,~$\star$~\nopagebreak

sie sollen vergehen vor deinem drohenden Angesicht.

\noindent Deine Hand schütze den Mann zu deiner Rechten,~$\star$~\nopagebreak

den Menschensohn, den du für dich groß und stark gemacht.

\noindent Erhalt uns am Leben!~$\star$~\nopagebreak

Dann wollen wir deinen Namen anrufen und nicht von dir weichen. 

\noindent Herr, Gott der Heerscharen, richte uns wieder auf!~$\star$~\nopagebreak

Lass dein Angesicht leuchten, dann ist uns geholfen.

\noindent Ehre sei dem Vater und dem Sohn~$\star$~\nopagebreak

und dem Heiligen Geist.

\noindent Wie im Anfang, so auch jetzt und alle Zeit~$\star$~\nopagebreak

und in Ewigkeit. Amen.

\vspace{10pt}

\noindent \textsc{Ant.} Sieh her, mein Gott, verbirg nicht dein Gesicht, denn mir ist angst, erhöre mich bald. 

\vspace{10pt}

\noindent \textsc{2 Ant.} Gott ist mein Retter; ihm will ich vertrauen und niemals verzagen. \nopagebreak

\subsubsection{Canticum}

\paragraph{Jes 12,1-6}

\noindent Ich danke dir, Herr.~†~\nopagebreak

Du hast mir gezürnt, doch dein Zorn hat sich gewendet~$\star$~\nopagebreak

und du hast mich getröstet.

\noindent Ja, Gott ist meine Rettung;~$\star$~\nopagebreak

ihm will ich vertrauen und niemals verzagen.

\noindent Denn meine Stärke und mein Lied ist der Herr.~$\star$~\nopagebreak

Er ist für mich zum Retter geworden.

\noindent Ihr werdet Wasser schöpfen voll Freude~$\star$~\nopagebreak

aus den Quellen des Heils.

\noindent An jenem Tag werdet ihr sagen:~$\star$~\nopagebreak

Dankt dem Herrn! Ruft seinen Namen an!

\noindent Macht seine Taten unter den Völkern bekannt,~$\star$~\nopagebreak

verkündet: Sein Name ist groß und erhaben!

\noindent Preist den Herrn, denn herrliche Taten hat er vollbracht;~$\star$~\nopagebreak

auf der ganzen Erde soll man es wissen.

\noindent Jauchzt und jubelt, ihr Bewohner von Zion,~$\star$~\nopagebreak

denn groß ist in eurer Mitte der Heilige Israels.

\noindent Ehre sei dem Vater und dem Sohn~$\star$~\nopagebreak

und dem Heiligen Geist.

\noindent Wie im Anfang, so auch jetzt und alle Zeit~$\star$~\nopagebreak

und in Ewigkeit. Amen.

\vspace{10pt}

\noindent \textsc{Ant.} Gott ist mein Retter; ihm will ich vertrauen und niemals verzagen. 

\vspace{10pt}

\noindent \textsc{3 Ant.} Mit bestem Weizen nährt uns der Herr und sättigt uns mit Honig aus dem Felsen.

\subsubsection{Psalm 81}

\noindent Jubelt Gott zu, er ist unsre Zuflucht;~$\star$~\nopagebreak

jauchzt dem Gott Jakobs zu!

\noindent Stimmt an den Gesang, schlagt die Pauke,~$\star$~\nopagebreak

die liebliche Laute, dazu die Harfe! 

\noindent Stoßt in die Posaune am Neumond~$\star$~\nopagebreak

und zum Vollmond, am Tag unsres Festes!

\noindent Denn das ist Satzung für Israel,~$\star$~\nopagebreak

Entscheid des Gottes Jakobs.

\noindent Das hat er als Gesetz für Josef erlassen,~$\star$~\nopagebreak

als Gott gegen Ägypten auszog.

\noindent Eine Stimme höre ich, die ich noch nie vernahm:~†~\nopagebreak

Seine Schulter hab ich von der Bürde befreit,~$\star$~\nopagebreak

seine Hände kamen los vom Lastkorb.

\noindent Du riefst in der Not~$\star$~\nopagebreak

und ich riss dich heraus;

\noindent ich habe dich aus dem Gewölk des Donners erhört,~$\star$~\nopagebreak

an den Wassern von Meríba geprüft.

\noindent Höre, mein Volk, ich will dich mahnen!~$\star$~\nopagebreak

Israel, wolltest du doch auf mich hören!

\noindent Für dich gibt es keinen andern Gott.~$\star$~\nopagebreak

Du sollst keinen fremden Gott anbeten.

\noindent Ich bin der Herr, dein Gott,~†~\nopagebreak

der dich heraufgeführt hat aus Ägypten.~$\star$~\nopagebreak

Tu deinen Mund auf! Ich will ihn füllen.

\noindent Doch mein Volk hat nicht auf meine Stimme gehört;~$\star$~\nopagebreak

Israel hat mich nicht gewollt.

\noindent Da überließ ich sie ihrem verstockten Herzen~$\star$~\nopagebreak

und sie handelten nach ihren eigenen Plänen. 

\noindent Ach dass doch mein Volk auf mich hörte,~$\star$~\nopagebreak

dass Israel gehen wollte auf meinen Wegen!

\noindent Wie bald würde ich seine Feinde beugen,~$\star$~\nopagebreak

meine Hand gegen seine Bedränger wenden. 

\noindent Alle, die den Herrn hassen, müssten Israel schmeicheln~$\star$~\nopagebreak

und das sollte für immer so bleiben. 

\noindent Ich würde es nähren mit bestem Weizen~$\star$~\nopagebreak

und mit Honig aus dem Felsen sättigen.

\noindent Ehre sei dem Vater und dem Sohn~$\star$~\nopagebreak

und dem Heiligen Geist.

\noindent Wie im Anfang, so auch jetzt und alle Zeit~$\star$~\nopagebreak

und in Ewigkeit. Amen.

\vspace{10pt}

\noindent \textsc{Ant.} Mit bestem Weizen nährt uns der Herr und sättigt uns mit Honig aus dem Felsen.
   \subsubsection{Kurzlesung}     \hfill Hebr 2, 9b--10        \lettrine[lines=3]{W}{}ir sehen Jesus um seines Todesleidens willen mit Herrlichkeit und Ehre gekrönt; es war nämlich Gottes gnädiger Wille, dass er für alle den Tod erlitt. Denn es war angemessen, dass Gott, für den und durch den das All ist und der viele Söhne zur Herrlichkeit führen wollte, den Urheber ihres Heils durch Leiden vollendete.
   \subsubsection{Responsorium}             \Rbar. Mit deinem heiligen Blute \GreStar{} hast du uns losgekauft.

\Rbar. Mit deinem heiligen Blute \GreStar{} hast du uns losgekauft.

\Vbar. Aus allen Stämmen und Sprachen. 

\Rbar. Hast du uns losgekauft.

\Vbar. Ehre sei dem Vater und dem Sohn und dem Heiligen Geist.

\Rbar. Mit deinem heiligen Blute \GreStar{} hast du uns losgekauft.
    \newpage
  \subsubsection{Benedictus} \noindent \textsc{Benedictus-Ant.} Mit Sehnsucht habe ich danach verlangt, dieses Ostermahl mit euch zu halten, bevor ich leide.

\subsubsection{Benedictus}

\paragraph{Lk 1, 68-79}

\gabcsnippet{(f3) (e f hr0 fR ,3) (hr0 iRr1 fR : hr0 e f gRr1 fR ::)}

\noindent Gepriesen sei der Herr, der Gott \uline{I}sraels!~$\star$~\nopagebreak

Denn er hat sein Volk besucht und ihm Erl\uline{ö}sung geschaffen;

\noindent er hat uns einen starken Retter erw\uline{e}ckt~$\star$~\nopagebreak

im Hause seines Kn\uline{e}chtes David.

\noindent So hat er verheißen von \uline{a}lters her~$\star$~\nopagebreak

durch den Mund seiner heilig\uline{e}n Propheten.

\noindent Er hat uns errettet vor unseren F\uline{ei}nden~$\star$~\nopagebreak

und aus der Hand aller, d\uline{ie} uns hassen;

\noindent er hat das Erbarmen mit den Vätern an uns vollendet~†~\nopagebreak

und an seinen heiligen B\uline{u}nd gedacht~$\star$~\nopagebreak

an den Eid, den er unserm Vater Abrah\uline{a}m geschworen hat;

\noindent er hat uns geschenkt, dass wir, aus Feindeshand befreit,~†~\nopagebreak

ihm furchtlos dienen in Heiligkeit und Ger\uline{e}chtigkeit~$\star$~\nopagebreak

vor seinem Angesicht all \uline{u}nsre Tage.

\noindent Und du, Kind, wirst Prophet des Höchst\uline{en hei}ßen;~†~\nopagebreak

denn du wirst dem Herrn vor\uline{a}ngehen~$\star$~\nopagebreak

und ihm den W\uline{e}g bereiten.

\noindent Du wirst sein Volk mit der Erfahrung des Heils b\uline{e}schenken~$\star$~\nopagebreak

in der Vergeb\uline{u}ng der Sünden. 

\noindent Durch die barmherzige Liebe unseres G\uline{o}ttes~$\star$~\nopagebreak

wird uns besuchen das aufstrahlende Licht \uline{au}s der Höhe,

\noindent um allen zu leuchten, die in Finsternis sitzen und im Schatten des T\uline{o}des,~$\star$~\nopagebreak

und unsre Schritte zu lenken auf den W\uline{e}g des Friedens.

\noindent Ehre sei dem Vater und dem S\uline{o}hn~$\star$~\nopagebreak

und dem H\uline{ei}ligen Geist.

\noindent Wie im Anfang, so auch jetzt und \uline{a}lle Zeit~$\star$~\nopagebreak

und in Ew\uline{i}gkeit. Amen.

\vspace{5pt}

\noindent \textsc{Ant.} Mit Sehnsucht habe ich danach verlangt, dieses Ostermahl mit euch zu halten, bevor ich leide.
   \subsubsection{Preces}  \greannotation{} \index[Preces]{Holy Thursday} \label{Holy Thursday (Preces)} \grecommentary[3pt]{} \gresetinitiallines{1} \grechangestyle{initial}{\fontsize{36}{36}\selectfont} \grechangedim{maxbaroffsettextleft@nobar}{12 cm}{scalable} \grechangedim{spaceabovelines}{0.7cm}{scalable} \gresetlyriccentering{vowel}   \gregorioscore{chants/misc.versus_litanici_in_cantu_feria_v}
   \subsubsection{Vater Unser}   \index[Vater Unser]{Vater Unser} \label{Vater Unser (Vater Unser)}             \vspace{10pt}
   \subsubsection{Oration}   \index[Oration]{Oration} \label{Oration (Oration)}         \lettrine[lines=3]{G}{}ott, es ist würdig und recht, dich über alles zu lieben. Mehre in uns den Reichtum deiner Gnade. Durch den Tod deines Sohnes lässt du uns erhoffen, was wir glauben. Gib, dass wir durch seine Auferstehung erlangen, was wir ersehnen. Darum bitten wir durch ihn, Jesus Christus.
\par \Rbar.~Amen.

   \subsubsection{Schlußsegen}   \index[Schlußsegen]{Schlußsegen} \label{Schlußsegen (Schlußsegen)}         Der Herr sei mit euch.

\Rbar. Und mit deinem Geiste.

Es segne euch der allmächtige Gott, + \par der Vater und der Sohn und der Heilige Geist.

\Rbar. Amen.

Gehet hin in Frieden.

\Rbar. Dank sei Gott, dem Herrn.


  \chapter{Trauermette am Karfreitag}
  \section{Officium Lectionis}
      \index[Varia]{Herr offne} \label{Herr offne (Varia)} \grecommentary[0pt]{} \gresetinitiallines{1} \grechangestyle{initial}{\fontsize{36}{36}\selectfont} \grechangedim{maxbaroffsettextleft@nobar}{12 cm}{scalable} \grechangedim{spaceabovelines}{0.5cm}{scalable} \gresetlyriccentering{syllable}  \grechangedim{maxbaroffsettextleft}{0 cm}{scalable} \gregorioscore{chants/herr_offne}
   \subsubsection{Invitatorium}   \index[Invitatorium]{Invitatorium} \label{Invitatorium (Invitatorium)}         \vspace{5pt} \par \gresetlyriccentering{syllable}

\emph{Cantor:}~
\nopagebreak
\greannotation{I}
\gabcsnippet{(c4) Im(d) Kreuz(h_) Je(g)su(h) Chri(j_)sti(h_) fin(g)den(f) wir(e_) Heil.(d_) (::)} \vspace{5pt}

\emph{Alle:}~
\nopagebreak
\gabcsnippet{(c4) Im(d) Kreuz(h_) Je(g)su(h) Chri(j_)sti(h_) fin(g)den(f) wir(e_) Heil.(d_) (::)} \vspace{5pt}

\emph{Cantor:}~
\nopagebreak
\gabcsnippet{Kommt(dh_), lasst(h) uns(h) ju(i)beln(h) vor(h) dem(g) Herrn(h) (,) und(j) zu(h)jauch(h)zen(g) dem(f) Fels(g) uns(g)res(h) Hei(f)les! (f) (:) 
Lasst(f) uns(f) mit(g) Lob(ixi_) sei(h)-nem(h) An(h)ge(h)sicht(f) na(gh)hen, (g) (,) vor(e) ihm(e) jau(g)chzen(g) mit(h) Lie(fe)dern! (d) <sp>R/</sp> (::)} \vspace{5pt}


\emph{Cantor:}~
\nopagebreak
\gabcsnippet{(c4)Denn(d) der(d) Herr(h_) ist(h) ein(h) gro(i)ßer(g) Gott,(h) (,) ein(h) gro(h)ßer(j) Kö(hi)nig(h) üb(h)er(h) all(g)en(h) Gö(f)ttern.(f) (:)  In(d) sei(f)ner(f) Hand(g_) sind(g) die(g) Tie(h)fen(g) der(h) Er(f)de,(g) (,) sein(f) sind(f) die(f) Gi(g)pfel(g) der(f) Ber(gh)ge.(h) (:) Sein(h) ist(g) das(h) Meer,(ixi_) das(h) er(h) ge(f)macht(g) hat,(g) (,) das(g) tro(h)cke(g)ne(f) Land,(g_) das(g) sei(g)ne(e) Hän(g)de(g)  ge(h)bil(fe)det.(d) (::)} \vspace{5pt}

\emph{Alle:}~
\nopagebreak
\gabcsnippet{(c4) Im(d) Kreuz(h_) Je(g)su(h) Chri(j_)sti(h_) fin(g)den(f) wir(e_) Heil.(d_) (::)} \vspace{5pt}

\emph{Cantor:}~
\nopagebreak

\gabcsnippet{(c4)Kommt(d_), lasst(d) uns(d) nie(h)der(h)fal(h)len(h), uns(h) vor(h) ihm(i) ver(g)nei(hi)gen,(h) (,) lasst(h) uns(j) nie(h)der(h)knien(h) vor(g) dem(f) Herrn,(g) un(g)serm(h) Schö(f)pfer!(f) (:) Denn(f) er(g) ist(g) un(f)ser(g) Gott,(h) (,) wir(h) sind(h) das(h) Volk(ixi_) sei(h)ner(f) Wei(gh)de,(g) (,) die(g) Her(fg)de,(g) von(e) sei(g)ner(h) Hand(f) ge(e)führt.(d) 
(::)}

 \vspace{5pt}


\emph{Alle:}~
\nopagebreak
\gabcsnippet{(c4) Im(d) Kreuz(h_) Je(g)su(h) Chri(j_)sti(h_) fin(g)den(f) wir(e_) Heil.(d_) (::)}

 \vspace{5pt}

\emph{Cantor:}~
\nopagebreak
\gabcsnippet{(c4) Ach,(d_) wür(d)det(d) ihr(d) doch(d) heu(dh)te(h) auf(h) sei(h)ne(h) Sti(i)mme(g) hö(hi)ren!(h) (;) „Ver(h)här(h)tet(h) eu(h)er(h) Herz(g) nicht(h) wie(h) in(h) Me(f)rí(g)ba,(g) (,) wie(h) in(h) der(j) Wü(hi)ste(h) am(h) Tag(g) von(h) Ma(f)ssa!(f) (:) Dort(f) ha(f)ben(f) eu(f)re(f) Vä(g)ter(g) mich(f) ver(g)sucht,(h) (,) sie(h) ha(h)ben(h) mich(h) auf(h) die(h) Pro(ixi)be(h) ge(f)stellt(g) und(g) hat(g)ten(g) doch(g) mein(e) Tun(g) ge(h)se(fe)hen.(d) 
(::)}

 \vspace{5pt}


\emph{Alle:}~
\nopagebreak
\gabcsnippet{(c4) Im(d) Kreuz(h_) Je(g)su(h) Chri(j_)sti(h_) fin(g)den(f) wir(e_) Heil.(d_) (::)}

 \vspace{5pt}

\emph{Cantor:}~\nopagebreak

\gabcsnippet{(c4) Vier(d)zig(d) Jah(dh)re(h) war(h) mir(h) dies(h) Ge(h)schlecht(i) zu(g)wi(hi)der,(h) (;) und(h) ich(f) sa(gh)gte:(g) Sie(h) sind(h) ein(j) Volk(h), de(h)ssen(h) Herz(h) in(g) die(f) Ir(g)re(g) geht;(g) (,) denn(h) mei(h)ne(j) We(hi)ge(h) ken(h)nen(g) sie(h) nicht.(f) (:) Da(f)rum(f) ha(f)be(f) ich(f) in(f) mei(f)nem(f) Zorn(g) ge(f)schwo(gh)ren:(h) (,) Sie(h) sol(ixi)len(h) nicht(f) kom(gh)men(g) in(g) das(e) Land(g) mei(g)ner(h) Ru(fe)he.”(d) (::)}

 \vspace{5pt}


\emph{Alle:}~\nopagebreak

\gabcsnippet{(c4) Im(d) Kreuz(h_) Je(g)su(h) Chri(j_)sti(h_) fin(g)den(f) wir(e_) Heil.(d_) (::)}

 \vspace{5pt}

\emph{Cantor:}~
\nopagebreak

\gabcsnippet{(c4) Eh(d)re(d) dem(d) Va(dh)ter,(h) (,) Eh(i)re(h) dem(g) Sohn,(h) (,) Eh(h)re(h) dem(j) Hei(h)li(g)gen(h) Geist.(f) (:) Wie(f) im(g) An(ixi)fang,(h) (,) so(h) auch(h) jetzt(h) und(f) al(g)le(h) Zeit(g) und(g) in(e) E(g)wig(g)keit.(h) A(fvED)men.(d) (::)}

 \vspace{5pt}


\emph{Alle:}~
\nopagebreak
\gabcsnippet{(c4) Im(d) Kreuz(h_) Je(g)su(h) Chri(j_)sti(h_) fin(g)den(f) wir(e_) Heil.(d_) (::)}
  \subsubsection{Hymnus}  \greannotation{I} \index[Hymnus]{Hymnus} \label{Hymnus (Hymnus)} \grecommentary[0pt]{} \gresetinitiallines{1} \grechangestyle{initial}{\fontsize{36}{36}\selectfont} \grechangedim{maxbaroffsettextleft@nobar}{12 cm}{scalable} \grechangedim{spaceabovelines}{0.5cm}{scalable} \gresetlyriccentering{syllable}   

\gabcsnippet{
(c4) Heil(d)ig(e) Kreuz,(gh) du(hg) Baum(h) der(j) Treu(ji)e,(g) (,)
ed(i)ler(j) Baum,(kl) dem(hg) kein(j)er(i) gleich,(h) (;)
kein(hj)er(i) so(g) an(e) Laub(f) und(dc) Blü(de)te,(e) (,)
kein(d)er(h) so(hg) an(ed) Früch(f)ten(e) reich:(d) (;)
Sü(hj)ßes(i) Holz,(g) o(e) sü(f)ße(dc) Nä(de)gel,(e) (,)
wel(d)che(h) sü(hg)ße(ed) Last(f) an(e) euch.(d) (::)}

\gresetinitiallines{0}

\vspace{7pt}

\gabcsnippet{2. Beu(d)ge,(e) ho(gh)her(hg) Baum,(h) die(j) Zwei(ji)ge,(g) (,)
wer(i)de(j) weich(kl) an(hg) Stamm(j) und(i) Ast,(h) (;)
denn(hj) dein(i) har(g)tes(e) Holz(f) muß(dc) tra(de)gen(e) (,)
ei(d)ne(h) kö(hg)nig(ed)li(f)che(e) Last,(d) (;)
gib(hj) den(i) Glie(g)dern(e) dei(f)nes(dc) Schöp(de)fers(e) (,)
an(d) dem(h) Stam(hg)me(ed) lin(f)de(e) Rast.(d) (::)}

\vspace{7pt}

\gabcsnippet{3. Du(d) al(e)lein(gh) warst(hg) wert,(h) zu(j) tra(ji)gen(g) (,)
al(i)ler(j) Sün(kl)den(hg) Lö(j)se(i)geld,(h) (;)
du,(hj) die(i) Plan(g)ke,(e) die(f) uns(dc) ret(de)tet(e) (,)
aus(d) dem(h) Schiff(hg)bruch(ed) die(f)ser(e) Welt.(d) (;)
Du,(hj) ge(i)salbt(g) vom(e) Blut(f) des(dc) Lam(de)mes,(e) (,)
Pfos(d)ten,(h) der(hg) den(ed) Tod(f) ab(e)hält.(d) (::)}

\vspace{7pt}

\gabcsnippet{4. Lob(d) und(e) Ruhm(gh) sei(hg) oh(h)ne(j) En(ji)de(g) (,)
Gott,(i) dem(j) höchs(kl)ten(hg) Herrn,(j) ge(i)weiht.(h) (;)
Preis(hj) dem(i) Va(g)ter(e) und(f) dem(dc) Soh(de)ne(e) (,)
und(d) dem(h) Geist(hg) der(ed) Hei(f)lig(e)keit.(d) (;)
Ei(hj)nen(i) Gott(g) in(e) drei(f) Per(dc)so(de)nen(e) (,)
lo(d)be(h) al(hg)le(ed) Welt(f) und(e) Zeit.(d) (::)
A(ded)men.(cd) (::)}                 \newpage
  \subsection{Psalmodie} \noindent \textsc{1 Ant.} Die Könige der Erde stehen auf, die Großen haben sich verbündet gegen den Herrn und seinen Gesalbten.

\subsubsection{Ps 2,1-12}

Warum toben die Völker,~$\star$~\nopagebreak

warum machen die Nationen vergebliche Pläne?
 
\noindent Die Könige der Erde stehen auf,~$\star$~\nopagebreak

die Großen haben sich verbündet gegen den Herrn und seinen Gesalbten.
 
\noindent «Lasst uns ihre Fesseln zerreißen~$\star$~\nopagebreak

und von uns werfen ihre Stricke!»
 
\noindent Doch er, der im Himmel thront, lacht,~$\star$~\nopagebreak

der Herr verspottet sie.
 
\noindent Dann aber spricht er zu ihnen im Zorn,~$\star$~\nopagebreak

in seinem Grimm wird er sie erschrecken:
 
\noindent «Ich selber habe meinen König eingesetzt~$\star$~\nopagebreak

auf Zion, meinem heiligen Berg.»
 
\noindent Den Beschluss des Herrn will ich kundtun.~†~\nopagebreak

Er sprach zu mir: «Mein Sohn bist du.~$\star$~\nopagebreak

Heute habe ich dich gezeugt.
 
\noindent Fordre von mir und ich gebe dir die Völker zum Erbe,~$\star$~\nopagebreak

die Enden der Erde zum Eigentum.
 
\noindent Du wirst sie zerschlagen mit eiserner Keule,~$\star$~\nopagebreak

wie Krüge aus Ton wirst du sie zertrümmern.»
 
\noindent Nun denn, ihr Könige, kommt zur Einsicht,~$\star$~\nopagebreak

lasst euch warnen, ihr Gebieter der Erde!
 
\noindent Dient dem Herrn in Furcht~$\star$~\nopagebreak

und küsst ihm mit Beben die Füße,
 
\noindent damit er nicht zürnt~$\star$~\nopagebreak

und euer Weg nicht in den Abgrund führt. 
 
\noindent Denn wenig nur und sein Zorn ist entbrannt.~$\star$~\nopagebreak

Wohl allen, die ihm vertrauen!

\noindent Ehre sei dem Vater und dem Sohn~$\star$~\nopagebreak

und dem Heiligen Geist.

\noindent Wie im Anfang, so auch jetzt und alle Zeit~$\star$~\nopagebreak

und in Ewigkeit. Amen.

\vspace{10pt}

\noindent \textsc{Ant.} Die Könige der Erde stehen auf, die Großen haben sich verbündet gegen den Herrn und seinen Gesalbten.

\vspace{10pt}


\noindent \textsc{2 Ant.} Sie verteilen unter sich meine Kleider und werfen das Los um mein Gewand.

\subsubsection{Ps 22 (21), 2-23}

Mein Gott, mein Gott, warum hast du mich verlassen,~$\star$~\nopagebreak

bist fern meinem Schreien, den Worten meiner Klage?
 
\noindent Mein Gott, ich rufe bei Tag, doch du gibst keine Antwort;~$\star$~\nopagebreak

ich rufe bei Nacht und finde doch keine Ruhe.
 
\noindent Aber du bist heilig,~$\star$~\nopagebreak

du thronst über dem Lobpreis Israels.
 
\noindent Dir haben unsre Väter vertraut,~$\star$~\nopagebreak

sie haben vertraut und du hast sie gerettet.
 
\noindent Zu dir riefen sie und wurden befreit,~$\star$~\nopagebreak

dir vertrauten sie und wurden nicht zuschanden.
 
\noindent Ich aber bin ein Wurm und kein Mensch,~$\star$~\nopagebreak

der Leute Spott, vom Volk verachtet.
 
\noindent Alle, die mich sehen, verlachen mich,~$\star$~\nopagebreak

verziehen die Lippen, schütteln den Kopf:
 
\noindent «Er wälze die Last auf den Herrn,~$\star$~\nopagebreak

der soll ihn befreien! 
 
\noindent Der reiße ihn heraus,~$\star$~\nopagebreak

wenn er an ihm Gefallen hat.»
 
\noindent Du bist es, der mich aus dem Schoß meiner Mutter zog,~$\star$~\nopagebreak

mich barg an der Brust der Mutter.
 
\noindent Von Geburt an bin ich geworfen auf dich,~$\star$~\nopagebreak

vom Mutterleib an bist du mein Gott.
 
\noindent Sei mir nicht fern, denn die Not ist nahe~$\star$~\nopagebreak

und niemand ist da, der hilft.
 
\noindent Viele Stiere umgeben mich,~$\star$~\nopagebreak

Büffel von Baschan umringen mich.
 
\noindent Sie sperren gegen mich ihren Rachen auf,~$\star$~\nopagebreak

reißende, brüllende Löwen.
 
\noindent Ich bin hingeschüttet wie Wasser,~†~\nopagebreak

gelöst haben sich all meine Glieder.~$\star$~\nopagebreak

Mein Herz ist in meinem Leib wie Wachs zerflossen.
 
\noindent Meine Kehle ist trocken wie eine Scherbe,~†~\nopagebreak

die Zunge klebt mir am Gaumen,~$\star$~\nopagebreak

du legst mich in den Staub des Todes.
 
\noindent Viele Hunde umlagern mich,~†~\nopagebreak

eine Rotte von Bösen umkreist mich.~$\star$~\nopagebreak

Sie durchbohren mir Hände und Füße.
 
\noindent Man kann all meine Knochen zählen;~$\star$~\nopagebreak

sie gaffen und weiden sich an mir.
 
\noindent Sie verteilen unter sich meine Kleider~$\star$~\nopagebreak

und werfen das Los um mein Gewand.
 
\noindent Du aber, Herr, halte dich nicht fern!~$\star$~\nopagebreak

Du, meine Stärke, eil mir zu Hilfe!
 
\noindent Entreiße mein Leben dem Schwert,~$\star$~\nopagebreak

mein einziges Gut aus der Gewalt der Hunde!
 
\noindent Rette mich vor dem Rachen des Löwen,~$\star$~\nopagebreak

vor den Hörnern der Büffel rette mich Armen!
 
\noindent Ich will deinen Namen meinen Brüdern verkünden,~$\star$~\nopagebreak

inmitten der Gemeinde dich preisen.

\noindent Ehre sei dem Vater und dem Sohn~$\star$~\nopagebreak

und dem Heiligen Geist.

\noindent Wie im Anfang, so auch jetzt und alle Zeit~$\star$~\nopagebreak

und in Ewigkeit. Amen.

\vspace{10pt}

\noindent \textsc{Ant.} Sie verteilen unter sich meine Kleider und werfen das Los um mein Gewand.

\vspace{10pt}

\noindent \textsc{3 Ant.} Die mir nach dem Leben trachten, legen mir Schlingen; die mein Unheil suchen, planen Verderben.

\subsubsection{Ps 38 (37)}

Herr, strafe mich nicht in deinem Zorn~$\star$~\nopagebreak

und züchtige mich nicht in deinem Grimm!
 
\noindent Denn deine Pfeile haben mich getroffen,~$\star$~\nopagebreak

deine Hand lastet schwer auf mir.
 
\noindent Nichts blieb gesund an meinem Leib, weil du mir grollst;~$\star$~\nopagebreak

weil ich gesündigt, blieb an meinen Gliedern nichts heil.
 
\noindent Denn meine Sünden schlagen mir über dem Kopf zusammen,~$\star$~\nopagebreak

sie erdrücken mich wie eine schwere Last.
 
\noindent Mir schwären, mir eitern die Wunden~$\star$~\nopagebreak

wegen meiner Torheit.
 
\noindent Ich bin gekrümmt und tief gebeugt,~$\star$~\nopagebreak

den ganzen Tag geh ich traurig einher.
 
\noindent Denn meine Lenden sind voller Brand,~$\star$~\nopagebreak

nichts blieb gesund an meinem Leib.
 
\noindent Kraftlos bin ich und ganz zerschlagen,~$\star$~\nopagebreak

ich schreie in der Qual meines Herzens.
 
\noindent All mein Sehnen, Herr, liegt offen vor dir,~$\star$~\nopagebreak

mein Seufzen ist dir nicht verborgen.
 
\noindent Mein Herz pocht heftig, mich hat die Kraft verlassen,~$\star$~\nopagebreak

geschwunden ist mir das Licht der Augen.
 
\noindent Freunde und Gefährten bleiben mir fern in meinem Unglück~$\star$~\nopagebreak

und meine Nächsten meiden mich.
 
\noindent Die mir nach dem Leben trachten, legen mir Schlingen;~†~\nopagebreak

die mein Unheil suchen, planen Verderben,~$\star$~\nopagebreak

den ganzen Tag haben sie Arglist im Sinn.
 
\noindent Ich bin wie ein Tauber, der nicht hört,~$\star$~\nopagebreak

wie ein Stummer, der den Mund nicht auftut.
 
\noindent Ich bin wie einer, der nicht mehr hören kann,~$\star$~\nopagebreak

aus dessen Mund keine Entgegnung kommt.
 
\noindent Doch auf dich, Herr, harre ich;~$\star$~\nopagebreak

du wirst mich erhören, Herr, mein Gott.
 
\noindent Denn ich sage: Über mich sollen die sich nicht freuen,~$\star$~\nopagebreak

die gegen mich prahlen, wenn meine Füße straucheln.
 
\noindent Ich bin dem Fallen nahe,~$\star$~\nopagebreak

mein Leid steht mir immer vor Augen.
 
\noindent Ja, ich bekenne meine Schuld,~$\star$~\nopagebreak

ich bin wegen meiner Sünde in Angst.
 
\noindent Die mich ohne Grund befehden, sind stark;~$\star$~\nopagebreak

viele hassen mich wegen nichts.
 
\noindent Sie vergelten mir Gutes mit Bösem,~$\star$~\nopagebreak

sie sind mir Feind; denn ich trachte nach dem Guten.
 
\noindent Herr, verlass mich nicht, bleib mir nicht fern, mein Gott!~$\star$~\nopagebreak

Eile mir zu Hilfe, Herr, du mein Heil!

\noindent Ehre sei dem Vater und dem Sohn~$\star$~\nopagebreak

und dem Heiligen Geist.

\noindent Wie im Anfang, so auch jetzt und alle Zeit~$\star$~\nopagebreak

und in Ewigkeit. Amen.

\vspace{10pt}

\noindent \textsc{Ant.} Die mir nach dem Leben trachten, legen mir Schlingen; die mein Unheil suchen, planen Verderben.


   \subsubsection{Versiculum}  \greannotation{} \index[Versiculum]{Versiculum} \label{Versiculum (Versiculum)} \grecommentary[0pt]{} \gresetinitiallines{0} \grechangestyle{initial}{\fontsize{36}{36}\selectfont} \grechangedim{maxbaroffsettextleft@nobar}{12 cm}{scalable} \grechangedim{spaceabovelines}{0.5cm}{scalable} \gresetlyriccentering{syllable}   \gregorioscore{chants/versiculum_thursday}
  \subsection{Lesungen}
   \subsubsection{Erste Lesung} \vspace{5pt} \emph{Aus den Klageliedern des Propheten Jeremia. Aleph. Weh, mit seinem Zorn umwölkt der Herr die Tochter Zion. Er schleudert vom Himmel zur Erde die Pracht Israels. Nicht dachte er an den Schemel seiner Füße am Tag seines Zornes. Beth. Schonungslos hat der Herr vernichtet alle Fluren Jakobs, niedergerissen in seinem Grimm die Bollwerke der Tochter Juda, zu Boden gestreckt, entweiht das Königtum und seine Fürsten. Ghimel. Abgehauen hat er in Zornesglut jedes Horn in Israel. Er zog seine Rechte zurück angesichts des Feindes und brannte in Jakob wie flammendes Feuer, ringsum alles verzehrend. - Jerusalem, Jerusalem, bekehre dich zum Herrn, deinem Gott.} \vspace{5pt} \greannotation{} \index[Erste Lesung]{Erste Lesung} \label{Erste Lesung (Erste Lesung)} \grecommentary[0pt]{Klgl 2, 1--3} \gresetinitiallines{1} \grechangestyle{initial}{\fontsize{36}{36}\selectfont} \grechangedim{maxbaroffsettextleft@nobar}{12 cm}{scalable} \grechangedim{spaceabovelines}{0.5cm}{scalable} \gresetlyriccentering{vowel}   \gregorioscore{chants/va--lamentationes_61_de_lamentatione_ieremiae_prophetae--dominican}    \newpage
   \subsubsection{II} \vspace{5pt} \emph{Daleth. Er spannte den Bogen wie ein Feind, stand da, erhoben die Rechte. Wie ein Gegner erschlug er alles, was das Auge erfreut. Im Zelt der Tochter Zion goss er seinen Zorn aus wie Feuer. He. Wie ein Feind ist geworden der Herr, Israel hat er vernichtet. Vernichtet hat er alle Paläste, zerstört seine Burgen. Auf die Tochter Juda hat er gehäuft Jammer über Jammer. Vau. Er zertrat wie einen Garten seine Wohnstatt, zerstörte seinen Festort. Vergessen ließ der Herr auf Zion Festtag und Sabbat. In glühendem Zorn verwarf er König und Priester. - Jerusalem, Jerusalem, bekehre dich zum Herrn, deinem Gott.} \vspace{5pt} \greannotation{} \index[II]{II} \label{II (II)} \grecommentary[0pt]{Klgl 2, 4--6} \gresetinitiallines{1} \grechangestyle{initial}{\fontsize{36}{36}\selectfont} \grechangedim{maxbaroffsettextleft@nobar}{12 cm}{scalable} \grechangedim{spaceabovelines}{0.5cm}{scalable} \gresetlyriccentering{vowel}   \gregorioscore{chants/va--lamentationes_62_daleth_tetendit--dominican}    \newpage
   \subsubsection{III} \vspace{5pt} \emph{Zain. Seinen Altar hat der Herr verschmäht, entweiht sein Heiligtum, überliefert in die Hand des Feindes die Mauern von Zions Palästen. Man lärmte im Haus des Herrn wie an einem Festtag. Zu schleifen plante der Herr die Mauer der Tochter Zion. Heth. Er spannte die Messschnur und zog nicht zurück die Hand vom Vertilgen. Trauern ließ er Wall und Mauer; miteinander sanken sie nieder. Teth. In den Boden sanken ihre Tore, ihre Riegel hat er zerstört und zerbrochen. Ihr König und ihre Fürsten sind unter den Völkern, keine Weisung ist da, auch keine Offenbarung schenkt der Herr ihren Propheten. - Jerusalem, Jerusalem, bekehre dich zum Herrn, deinem Gott.}  \vspace{5pt} \greannotation{} \index[III]{III} \label{III (III)} \grecommentary[0pt]{Klgl 2, 7--9} \gresetinitiallines{1} \grechangestyle{initial}{\fontsize{36}{36}\selectfont} \grechangedim{maxbaroffsettextleft@nobar}{12 cm}{scalable} \grechangedim{spaceabovelines}{0.5cm}{scalable} \gresetlyriccentering{vowel}   \gregorioscore{chants/va--lamentationes_63_zain_repulit--dominican}    \newpage
   \subsubsection{Zweite Lesung}             \vspace{10pt}

Johannes Chrysostomus († 407), Aus einer Katechese


\vspace{10pt}



\lettrine[lines=3]{W}{}illst du erfahren, welche Kraft das Blut Christi besitzt? Dann laß uns zurückgehen bis zu dem Vorausbild. Auf das frühe Vorausbild wollen wir uns besinnen und die Niederschrift aus der Vergangenheit erzählen. Mose sagt: „Tötet ein einjähriges Lamm und bestreicht mit seinem Blut die Tür.“ Was sagst du da, Mose? Kann denn das Blut eines Lammes den vernunftbegabten Menschen befreien? Gewiß, sagt er, weil es auf das Blut des Herrn verweist. Wenn der Feind nicht das Blut des Vorbildes an Pfosten, sondern auf den Lippen der Glaubenden das kostbare Blut der Wahrheit leuchten sieht, mit dem der Tempel Christi geweiht ist, dann weicht er viel weiter zurück. 

Willst du der Kraft dieses Blutes noch weiter nachforschen? Dann schau bitte, woher es kommt und aus welcher Quelle es entspringt. Vom Kreuz Christi kam es zuerst, aus der Seite Christi nahm es den Anfang. Denn das Evangelium berichtet: Als Jesus tot war und noch am Kreuz hing, kam ein Soldat herbei und stieß die Seite auf. Da floß Wasser und Blut heraus: Symbol der Taufe das eine, Symbol des Mysteriums (der Eucharistie) das andere. Der Soldat hat die Seite geöffnet und die Wand des Tempels aufgetan. Ich habe den herrlichen Schatz gefunden und bin glücklich, den glanzvollen Reichtum entdeckt zu haben. So war es auch mit dem Lamm: Die Juden haben es geschlachtet, und ich erfahre die Frucht des Opfers. 

Blut und Wasser aus der Seite. Lieber Hörer, bitte geh nicht eilig an dem verborgenen Mysterium vorbei. Denn ich muß noch mystische und geheime Dinge aussprechen: Ich sagte, dieses Wasser und Blut seien Sinnzeichen für die Taufe und das Mysterium. Daraus ist die heilige Kirche aufgebaut, durch die Wiedergeburt aus dem Wasser und die Erneuerung des Heiligen Geistes, ich sage euch: durch die Taufe und das Mysterium, das aus seiner Seite hervorging. Aus seiner Seite nämlich baute Christus die Kirche, wie aus der Seite Adams Eva, die Gattin, kam. 

Dafür ist auch Paulus Zeuge, wenn er sagt: „Wir sind Glieder seines Leibes“, von seinem Gebein genommen, womit er die Seite meint. Denn wie Gott aus der Seite des Adam die Frau schuf, so gab uns Christus aus seiner Seite Wasser und Blut, wodurch die Kirche erbaut werden sollte. Wie Gott die Seite öffnete, während Adam im Schlaf ruhte, so schenkte er uns jetzt nach dem Tode Christi aus seiner Seite das Wasser und das Blut.
    \newpage
  \section{Laudes}
  \subsubsection{Psalmodie} \noindent \textsc{1 Ant.} Seinen eigenen Sohn hat Gott nicht verschont: Er hat ihn hingegeben für uns alle.

\subsubsection{Psalm 51 (50), 3--21}

\noindent Gott, sei mir gnädig nach deiner Huld,~\GreStar{}~\nopagebreak

tilge meine Frevel nach deinem reichen Erbarmen!

\noindent Wasch meine Schuld von mir ab~\GreStar{}~\nopagebreak

und mach mich rein von meiner Sünde!

\noindent Denn ich erkenne meine bösen Taten,~\GreStar{}~\nopagebreak

meine Sünde steht mir immer vor Augen.

\noindent Gegen dich allein habe ich gesündigt,~\GreStar{}~\nopagebreak

ich habe getan, was dir missfällt. 

\noindent So behältst du recht mit deinem Urteil,~\GreStar{}~\nopagebreak

rein stehst du da als Richter.

\noindent Denn ich bin in Schuld geboren;~\GreStar{}~\nopagebreak

in Sünde hat mich meine Mutter empfangen.

\noindent Lauterer Sinn im Verborgenen gefällt dir,~\GreStar{}~\nopagebreak

im Geheimen lehrst du mich Weisheit.

\noindent Entsündige mich mit Ysop, dann werde ich rein;~\GreStar{}~\nopagebreak

wasche mich, dann werde ich weißer als Schnee.

\noindent Sättige mich mit Entzücken und Freude!~\GreStar{}~\nopagebreak

Jubeln sollen die Glieder, die du zerschlagen hast.

\noindent Verbirg dein Gesicht vor meinen Sünden;~\GreStar{}~\nopagebreak

tilge all meine Frevel!

\noindent Erschaffe mir, Gott, ein reines Herz~\GreStar{}~\nopagebreak

und gib mir einen neuen, beständigen Geist!

\noindent Verwirf mich nicht von deinem Angesicht~\GreStar{}~\nopagebreak

und nimm deinen heiligen Geist nicht von mir!

\noindent Mach mich wieder froh mit deinem Heil;~\GreStar{}~\nopagebreak

mit einem willigen Geist rüste mich aus!

\noindent Dann lehre ich Abtrünnige deine Wege,~\GreStar{}~\nopagebreak

und die Sünder kehren um zu dir.

\noindent Befrei mich von Blutschuld, Herr, du Gott meines Heiles,~\GreStar{}~\nopagebreak

dann wird meine Zunge jubeln über deine Gerechtigkeit.

\noindent Herr, öffne mir die Lippen,~\GreStar{}~\nopagebreak

und mein Mund wird deinen Ruhm verkünden.

\noindent Schlachtopfer willst du nicht, ich würde sie dir geben;~\GreStar{}~\nopagebreak

an Brandopfern hast du kein Gefallen.

\noindent Das Opfer, das Gott gefällt, ist ein zerknirschter Geist,~\GreStar{}~\nopagebreak

ein zerbrochenes und zerschlagenes Herz 

wirst du, Gott, nicht verschmähen.

\noindent In deiner Huld tu Gutes an Zion;~\GreStar{}~\nopagebreak

bau die Mauern Jerusalems wieder auf!

\noindent Dann hast du Freude an rechten Opfern,~†~\nopagebreak

an Brandopfern und Ganzopfern,~\GreStar{}~\nopagebreak

dann opfert man Stiere auf deinem Altar.

\noindent Ehre sei dem Vater und dem Sohn~\GreStar{}~\nopagebreak

und dem Heiligen Geist.

\noindent Wie im Anfang, so auch jetzt und alle Zeit~\GreStar{}~\nopagebreak

und in Ewigkeit. Amen.

\vspace{10pt}

\noindent \textsc{Ant.} Seinen eigenen Sohn hat Gott nicht verschont: Er hat ihn hingegeben für uns alle.

\vspace{10pt}

\noindent \textsc{2 Ant.} Jesus Christus hat uns geliebt und durch sein Blut von unseren Sünden befreit. \nopagebreak

\subsubsection{Canticum}

\paragraph{Hab 3, 2--4.13a.15--19}

\noindent Herr, ich höre die Kunde,~\GreStar{}~\nopagebreak

ich sehe, Herr, was du früher getan hast. 

\noindent Lass es in diesen Jahren wieder geschehen,~\GreStar{}~\nopagebreak

offenbare es in diesen Jahren! 

\noindent Auch wenn du zürnst,~\GreStar{}~\nopagebreak

denk an dein Erbarmen!

\noindent Gott kommt von Teman her,~\GreStar{}~\nopagebreak

der Heilige kommt vom Gebirge Paran. 

\noindent Seine Hoheit überstrahlt den Himmel,~\GreStar{}~\nopagebreak

sein Ruhm erfüllt die Erde.

\noindent Er leuchtet wie das Licht der Sonne,~†~\nopagebreak

ein Kranz von Strahlen umgibt ihn,~\GreStar{}~\nopagebreak

in ihnen verbirgt sich seine Macht.

\noindent Du ziehst aus, um dein Volk zu retten,~\GreStar{}~\nopagebreak

um deinem Gesalbten zu helfen.

\noindent Du bahnst mit deinen Rossen den Weg durch das Meer,~\GreStar{}~\nopagebreak

durch das gewaltig schäumende Wasser.

\noindent Ich zitterte am ganzen Leib, als ich es hörte,~\GreStar{}~\nopagebreak

ich vernahm den Lärm, und ich schrie. 

\noindent Fäulnis befällt meine Glieder,~\GreStar{}~\nopagebreak

und es wanken meine Schritte. 

\noindent Doch in Ruhe erwarte ich den Tag der Not,~\GreStar{}~\nopagebreak

der dem Volk bevorsteht, das über uns herfällt.

\noindent Zwar blüht der Feigenbaum nicht,~\GreStar{}~\nopagebreak

an den Reben ist nichts zu ernten, 

\noindent der Ölbaum bringt keinen Ertrag,~\GreStar{}~\nopagebreak

die Kornfelder tragen keine Frucht; 

\noindent im Pferch sind keine Schafe,~\GreStar{}~\nopagebreak

im Stall steht kein Rind mehr.

\noindent Dennoch will ich jubeln über den Herrn~\GreStar{}~\nopagebreak

und mich freuen über Gott, meinen Retter.

\noindent Gott, der Herr, ist meine Kraft.~†~\nopagebreak

Er macht meine Füße schnell wie die Füße der Hirsche~\GreStar{}~\nopagebreak

und lässt mich schreiten auf den Höhen.

\noindent Ehre sei dem Vater und dem Sohn~\GreStar{}~\nopagebreak

und dem Heiligen Geist.

\noindent Wie im Anfang, so auch jetzt und alle Zeit~\GreStar{}~\nopagebreak

und in Ewigkeit. Amen.

\vspace{10pt}

\noindent \textsc{Ant.} Jesus Christus hat uns geliebt und durch sein Blut von unseren Sünden befreit.

\vspace{10pt}

\noindent \textsc{3 Ant.} Dein Kreuz, o Herr, verehren wir, und deine heilige Auferstehung preisen und rühmen wir; denn siehe, durch das Holz des Kreuzes kam Freude in alle Welt.

\subsubsection{Psalm 47,12--20}

\noindent Jerusalem, preise den Herrn,~\GreStar{}~\nopagebreak

lobsinge, Zion, deinem Gott!

\noindent Denn er hat die Riegel deiner Tore festgemacht,~\GreStar{}~\nopagebreak

die Kinder in deiner Mitte gesegnet;

\noindent er verschafft deinen Grenzen Frieden~\GreStar{}~\nopagebreak

und sättigt dich mit bestem Weizen.

\noindent Er sendet sein Wort zur Erde,~\GreStar{}~\nopagebreak

rasch eilt sein Befehl dahin.

\noindent Er spendet Schnee wie Wolle,~\GreStar{}~\nopagebreak

streut den Reif aus wie Asche.

\noindent Eis wirft er herab in Brocken,~\GreStar{}~\nopagebreak

vor seiner Kälte erstarren die Wasser.

\noindent Er sendet sein Wort aus, und sie schmelzen,~\GreStar{}~\nopagebreak

er lässt den Wind wehen, dann rieseln die Wasser.

\noindent Er verkündet Jakob sein Wort,~\GreStar{}~\nopagebreak

Israel seine Gesetze und Rechte.

\noindent An keinem andern Volk hat er so gehandelt,~\GreStar{}~\nopagebreak

keinem sonst seine Rechte verkündet.

\noindent Ehre sei dem Vater und dem Sohn~\GreStar{}~\nopagebreak

und dem Heiligen Geist.

\noindent Wie im Anfang, so auch jetzt und alle Zeit~\GreStar{}~\nopagebreak

und in Ewigkeit. Amen.

\vspace{10pt}

\noindent \textsc{Ant.} Dein Kreuz, o Herr, verehren wir, und deine heilige Auferstehung preisen und rühmen wir; denn siehe, durch das Holz des Kreuzes kam Freude in alle Welt.
                 \newpage
   \subsubsection{Kurzlesung}     \hfill Jes 52,13--15        \lettrine[lines=3]{S}{}eht, mein Knecht hat Erfolg, er wird groß sein und hoch erhaben. Viele haben sich über ihn entsetzt, so entstellt sah er aus, nicht mehr wie ein Mensch, seine Gestalt war nicht mehr die eines Menschen. Jetzt aber setzt er viele Völker in Staunen, Könige müssen vor ihm verstummen. Denn was man ihnen noch nie erzählt hat, das sehen sie nun; was sie niemals hörten, das erfahren sie jetzt.
   \subsubsection{Responsorium} \emph{Christus war für uns gehorsam bis zum Tod, bis zum Tod am Kreuze.} \vspace{5pt}  \greannotation{V} \index[Responsorium]{Responsorium} \label{Responsorium (Responsorium)} \grecommentary[7pt]{Phil 2, 8} \gresetinitiallines{1} \grechangestyle{initial}{\fontsize{36}{36}\selectfont} \grechangedim{maxbaroffsettextleft@nobar}{12 cm}{scalable} \grechangedim{spaceabovelines}{0.5cm}{scalable} \gresetlyriccentering{vowel}   \gregorioscore{chants/gr--christus_factus_est--dominican--id_6737--only-respond}    \newpage
  \subsubsection{Benedictus} \noindent \textsc{Ant.} Über seinem Haupt befestigten sie eine Inschrift: Jesus von Nazaret, der König der Juden.

\paragraph{Lk 1, 68-79}

\gabcsnippet{(f3) (e f hr0 fR ,3) (hr0 iRr1 fR : hr0 e f gRr1 fR ::)}

\noindent Gepriesen sei der Herr, der Gott \uline{I}sraels!~$\star$~\nopagebreak

Denn er hat sein Volk besucht und ihm Erlös\uline{u}ng geschaffen;

\noindent er hat uns einen starken Retter erw\uline{e}ckt~$\star$~\nopagebreak

im Hause seines Kn\uline{e}chtes David.

\noindent So hat er verheißen von \uline{a}lters her~$\star$~\nopagebreak

durch den Mund seiner heilig\uline{e}n Propheten.

\noindent Er hat uns errettet vor unseren F\uline{ei}nden~$\star$~\nopagebreak

und aus der Hand aller, d\uline{ie} uns hassen;

\noindent er hat das Erbarmen mit den Vätern an uns vollend\uline{e}t~†~\nopagebreak

und an seinen heiligen B\uline{u}nd gedacht~$\star$~\nopagebreak

an den Eid, den er unserm Vater Abrah\uline{a}m geschworen hat;

\noindent er hat uns geschenkt, dass wir, aus Feindeshand befr\uline{ei}t,~†~\nopagebreak

ihm furchtlos dienen in Heiligkeit und Ger\uline{e}chtigkeit~$\star$~\nopagebreak

vor seinem Angesicht all \uline{u}nsre Tage.

\noindent Und du, Kind, wirst Prophet des Höchst\uline{en hei}ßen;~†~\nopagebreak

denn du wirst dem Herrn vor\uline{a}ngehn~$\star$~\nopagebreak

und ihm den W\uline{e}g bereiten.

\noindent Du wirst sein Volk mit der Erfahrung des Heils besch\uline{e}nken~$\star$~\nopagebreak

in der Vergeb\uline{u}ng der Sünden. 

\noindent Durch die barmherzige Liebe unseres G\uline{o}ttes~$\star$~\nopagebreak

wird uns besuchen das aufstrahlende Licht \uline{au}s der Höhe,

\noindent um allen zu leuchten, die in Finsternis sitzen und im Schatten des T\uline{o}des,~$\star$~\nopagebreak

und unsre Schritte zu lenken auf den W\uline{e}g des Friedens.

\noindent Ehre sei dem Vater und dem S\uline{o}hn~$\star$~\nopagebreak

und dem H\uline{ei}ligen Geist.

\noindent Wie im Anfang, so auch jetzt und \uline{a}lle Zeit~$\star$~\nopagebreak

und in Ew\uline{i}gkeit. Amen.

\vspace{5pt}

\noindent \textsc{Ant.} Über seinem Haupt befestigten sie eine Inschrift: Jesus von Nazaret, der König der Juden.                 \newpage
   \subsubsection{Preces}  \greannotation{} \index[Preces]{Holy Thursday} \label{Holy Thursday (Preces)} \grecommentary[3pt]{} \gresetinitiallines{1} \grechangestyle{initial}{\fontsize{36}{36}\selectfont} \grechangedim{maxbaroffsettextleft@nobar}{12 cm}{scalable} \grechangedim{spaceabovelines}{0.7cm}{scalable} \gresetlyriccentering{vowel}   \gregorioscore{chants/misc.versus_litanici_in_cantu_feria_vi--dominican-de-rubrics}
   \subsubsection{Vater Unser}   \index[Vater Unser]{Vater Unser} \label{Vater Unser (Vater Unser)}             \vspace{10pt}
   \subsubsection{Oration}   \index[Oration]{Oration} \label{Oration (Oration)}         \lettrine[lines=3]{H}{}err, unser Gott, sieh herab auf deine Familie, für die unser Herr Jesus Christus sich willig den Händen der Frevler überliefert und die Marter des Kreuzes auf sich genommen hat. Er, der in der Einheit des Heiligen Geistes mit dir lebt und herrscht in alle Ewigkeit.
\par \Rbar.~Amen.

   \subsubsection{Schlußsegen}   \index[Schlußsegen]{Schlußsegen} \label{Schlußsegen (Schlußsegen)}         Der Herr sei mit euch.

\Rbar. Und mit deinem Geiste.

Es segne euch der allmächtige Gott, + \par der Vater und der Sohn und der Heilige Geist.

\Rbar. Amen.

Gehet hin in Frieden.

\Rbar. Dank sei Gott, dem Herrn.









  \chapter{Trauermette am Karsamstag}
  \section{Officium Lectionis}
      \index[Varia]{Herr offne} \label{Herr offne (Varia)} \grecommentary[0pt]{} \gresetinitiallines{1} \grechangestyle{initial}{\fontsize{36}{36}\selectfont} \grechangedim{maxbaroffsettextleft@nobar}{12 cm}{scalable} \grechangedim{spaceabovelines}{0.5cm}{scalable} \gresetlyriccentering{syllable}  \grechangedim{maxbaroffsettextleft}{0 cm}{scalable} \gregorioscore{chants/herr_offne}
   \subsubsection{Invitatorium}   \index[Invitatorium]{Invitatorium} \label{Invitatorium (Invitatorium)}         \vspace{5pt} \par \gresetlyriccentering{syllable}

\emph{Cantor:}~
\nopagebreak
\greannotation{I}
\gabcsnippet{(c4) Im(d) Kreuz(h_) Je(g)su(h) Chri(j_)sti(h_) fin(g)den(f) wir(e_) Heil.(d_) (::)} \vspace{5pt}

\emph{Alle:}~
\nopagebreak
\gabcsnippet{(c4) Im(d) Kreuz(h_) Je(g)su(h) Chri(j_)sti(h_) fin(g)den(f) wir(e_) Heil.(d_) (::)} \vspace{5pt}

\emph{Cantor:}~
\nopagebreak
\gabcsnippet{Kommt(dh_), lasst(h) uns(h) ju(i)beln(h) vor(h) dem(g) Herrn(h) (,) und(j) zu(h)jauch(h)zen(g) dem(f) Fels(g) uns(g)res(h) Hei(f)les! (f) (:) 
Lasst(f) uns(f) mit(g) Lob(ixi_) sei(h)-nem(h) An(h)ge(h)sicht(f) na(gh)hen, (g) (,) vor(e) ihm(e) jau(g)chzen(g) mit(h) Lie(fe)dern! (d) <sp>R/</sp> (::)} \vspace{5pt}


\emph{Cantor:}~
\nopagebreak
\gabcsnippet{(c4)Denn(d) der(d) Herr(h_) ist(h) ein(h) gro(i)ßer(g) Gott,(h) (,) ein(h) gro(h)ßer(j) Kö(hi)nig(h) üb(h)er(h) all(g)en(h) Gö(f)ttern.(f) (:)  In(d) sei(f)ner(f) Hand(g_) sind(g) die(g) Tie(h)fen(g) der(h) Er(f)de,(g) (,) sein(f) sind(f) die(f) Gi(g)pfel(g) der(f) Ber(gh)ge.(h) (:) Sein(h) ist(g) das(h) Meer,(ixi_) das(h) er(h) ge(f)macht(g) hat,(g) (,) das(g) tro(h)cke(g)ne(f) Land,(g_) das(g) sei(g)ne(e) Hän(g)de(g)  ge(h)bil(fe)det.(d) (::)} \vspace{5pt}

\emph{Alle:}~
\nopagebreak
\gabcsnippet{(c4) Im(d) Kreuz(h_) Je(g)su(h) Chri(j_)sti(h_) fin(g)den(f) wir(e_) Heil.(d_) (::)} \vspace{5pt}

\emph{Cantor:}~
\nopagebreak

\gabcsnippet{(c4)Kommt(d_), lasst(d) uns(d) nie(h)der(h)fal(h)len(h), uns(h) vor(h) ihm(i) ver(g)nei(hi)gen,(h) (,) lasst(h) uns(j) nie(h)der(h)knien(h) vor(g) dem(f) Herrn,(g) un(g)serm(h) Schö(f)pfer!(f) (:) Denn(f) er(g) ist(g) un(f)ser(g) Gott,(h) (,) wir(h) sind(h) das(h) Volk(ixi_) sei(h)ner(f) Wei(gh)de,(g) (,) die(g) Her(fg)de,(g) von(e) sei(g)ner(h) Hand(f) ge(e)führt.(d) 
(::)}

 \vspace{5pt}


\emph{Alle:}~
\nopagebreak
\gabcsnippet{(c4) Im(d) Kreuz(h_) Je(g)su(h) Chri(j_)sti(h_) fin(g)den(f) wir(e_) Heil.(d_) (::)}

 \vspace{5pt}

\emph{Cantor:}~
\nopagebreak
\gabcsnippet{(c4) Ach,(d_) wür(d)det(d) ihr(d) doch(d) heu(dh)te(h) auf(h) sei(h)ne(h) Sti(i)mme(g) hö(hi)ren!(h) (;) „Ver(h)här(h)tet(h) eu(h)er(h) Herz(g) nicht(h) wie(h) in(h) Me(f)rí(g)ba,(g) (,) wie(h) in(h) der(j) Wü(hi)ste(h) am(h) Tag(g) von(h) Ma(f)ssa!(f) (:) Dort(f) ha(f)ben(f) eu(f)re(f) Vä(g)ter(g) mich(f) ver(g)sucht,(h) (,) sie(h) ha(h)ben(h) mich(h) auf(h) die(h) Pro(ixi)be(h) ge(f)stellt(g) und(g) hat(g)ten(g) doch(g) mein(e) Tun(g) ge(h)se(fe)hen.(d) 
(::)}

 \vspace{5pt}


\emph{Alle:}~
\nopagebreak
\gabcsnippet{(c4) Im(d) Kreuz(h_) Je(g)su(h) Chri(j_)sti(h_) fin(g)den(f) wir(e_) Heil.(d_) (::)}

 \vspace{5pt}

\emph{Cantor:}~\nopagebreak

\gabcsnippet{(c4) Vier(d)zig(d) Jah(dh)re(h) war(h) mir(h) dies(h) Ge(h)schlecht(i) zu(g)wi(hi)der,(h) (;) und(h) ich(f) sa(gh)gte:(g) Sie(h) sind(h) ein(j) Volk(h), de(h)ssen(h) Herz(h) in(g) die(f) Ir(g)re(g) geht;(g) (,) denn(h) mei(h)ne(j) We(hi)ge(h) ken(h)nen(g) sie(h) nicht.(f) (:) Da(f)rum(f) ha(f)be(f) ich(f) in(f) mei(f)nem(f) Zorn(g) ge(f)schwo(gh)ren:(h) (,) Sie(h) sol(ixi)len(h) nicht(f) kom(gh)men(g) in(g) das(e) Land(g) mei(g)ner(h) Ru(fe)he.”(d) (::)}

 \vspace{5pt}


\emph{Alle:}~\nopagebreak

\gabcsnippet{(c4) Im(d) Kreuz(h_) Je(g)su(h) Chri(j_)sti(h_) fin(g)den(f) wir(e_) Heil.(d_) (::)}

 \vspace{5pt}

\emph{Cantor:}~
\nopagebreak

\gabcsnippet{(c4) Eh(d)re(d) dem(d) Va(dh)ter,(h) (,) Eh(i)re(h) dem(g) Sohn,(h) (,) Eh(h)re(h) dem(j) Hei(h)li(g)gen(h) Geist.(f) (:) Wie(f) im(g) An(ixi)fang,(h) (,) so(h) auch(h) jetzt(h) und(f) al(g)le(h) Zeit(g) und(g) in(e) E(g)wig(g)keit.(h) A(fvED)men.(d) (::)}

 \vspace{5pt}


\emph{Alle:}~
\nopagebreak
\gabcsnippet{(c4) Im(d) Kreuz(h_) Je(g)su(h) Chri(j_)sti(h_) fin(g)den(f) wir(e_) Heil.(d_) (::)}
  \subsubsection{Hymnus}  \greannotation{I} \index[Hymnus]{Hymnus} \label{Hymnus (Hymnus)} \grecommentary[0pt]{} \gresetinitiallines{1} \grechangestyle{initial}{\fontsize{36}{36}\selectfont} \grechangedim{maxbaroffsettextleft@nobar}{12 cm}{scalable} \grechangedim{spaceabovelines}{0.5cm}{scalable} \gresetlyriccentering{syllable}   

\gabcsnippet{
(c4) Heil(d)ig(e) Kreuz,(gh) du(hg) Baum(h) der(j) Treu(ji)e,(g) (,)
ed(i)ler(j) Baum,(kl) dem(hg) kein(j)er(i) gleich,(h) (;)
kein(hj)er(i) so(g) an(e) Laub(f) und(dc) Blü(de)te,(e) (,)
kein(d)er(h) so(hg) an(ed) Früch(f)ten(e) reich:(d) (;)
Sü(hj)ßes(i) Holz,(g) o(e) sü(f)ße(dc) Nä(de)gel,(e) (,)
wel(d)che(h) sü(hg)ße(ed) Last(f) an(e) euch.(d) (::)}

\gresetinitiallines{0}

\vspace{7pt}

\gabcsnippet{2. Beu(d)ge,(e) ho(gh)her(hg) Baum,(h) die(j) Zwei(ji)ge,(g) (,)
wer(i)de(j) weich(kl) an(hg) Stamm(j) und(i) Ast,(h) (;)
denn(hj) dein(i) har(g)tes(e) Holz(f) muß(dc) tra(de)gen(e) (,)
ei(d)ne(h) kö(hg)nig(ed)li(f)che(e) Last,(d) (;)
gib(hj) den(i) Glie(g)dern(e) dei(f)nes(dc) Schöp(de)fers(e) (,)
an(d) dem(h) Stam(hg)me(ed) lin(f)de(e) Rast.(d) (::)}

\vspace{7pt}

\gabcsnippet{3. Du(d) al(e)lein(gh) warst(hg) wert,(h) zu(j) tra(ji)gen(g) (,)
al(i)ler(j) Sün(kl)den(hg) Lö(j)se(i)geld,(h) (;)
du,(hj) die(i) Plan(g)ke,(e) die(f) uns(dc) ret(de)tet(e) (,)
aus(d) dem(h) Schiff(hg)bruch(ed) die(f)ser(e) Welt.(d) (;)
Du,(hj) ge(i)salbt(g) vom(e) Blut(f) des(dc) Lam(de)mes,(e) (,)
Pfos(d)ten,(h) der(hg) den(ed) Tod(f) ab(e)hält.(d) (::)}

\vspace{7pt}

\gabcsnippet{4. Lob(d) und(e) Ruhm(gh) sei(hg) oh(h)ne(j) En(ji)de(g) (,)
Gott,(i) dem(j) höchs(kl)ten(hg) Herrn,(j) ge(i)weiht.(h) (;)
Preis(hj) dem(i) Va(g)ter(e) und(f) dem(dc) Soh(de)ne(e) (,)
und(d) dem(h) Geist(hg) der(ed) Hei(f)lig(e)keit.(d) (;)
Ei(hj)nen(i) Gott(g) in(e) drei(f) Per(dc)so(de)nen(e) (,)
lo(d)be(h) al(hg)le(ed) Welt(f) und(e) Zeit.(d) (::)
A(ded)men.(cd) (::)}                 \newpage
  \subsection{Psalmodie} \noindent \textsc{1 Ant.} Ich lege mich nieder und ruhe in Frieden.

\subsubsection{Psalm 4, 2--9}

\noindent Wenn ich rufe, erhöre mich,~\GreStar{}~\nopagebreak

Gott, du mein Retter! 

\noindent Du hast mir Raum geschaffen, als mir angst war.~\GreStar{}~\nopagebreak

Sei mir gnädig und hör auf mein Flehen!

\noindent Ihr Mächtigen, wie lange noch schmäht ihr meine Ehre,~\GreStar{}~\nopagebreak

warum liebt ihr den Schein und sinnt auf Lügen?

\noindent Erkennt doch: Wunderbar handelt der Herr an den Frommen;~\GreStar{}~\nopagebreak

der Herr erhört mich, wenn ich zu ihm rufe.

\noindent Ereifert ihr euch, so sündigt nicht!~\GreStar{}~\nopagebreak

Bedenkt es auf eurem Lager und werdet stille!

\noindent Bringt rechte Opfer dar~\GreStar{}~\nopagebreak

und vertraut auf den Herrn!

\noindent Viele sagen: „Wer lässt uns Gutes erleben?“~\GreStar{}~\nopagebreak

Herr, lass dein Angesicht über uns leuchten!

\noindent Du legst mir größere Freude ins Herz,~\GreStar{}~\nopagebreak

als andere haben bei Korn und Wein in Fülle.

\noindent In Frieden leg’ ich mich nieder und schlafe ein;~\GreStar{}~\nopagebreak

denn du allein, Herr, lässt mich sorglos ruhen.

\noindent Ehre sei dem Vater und dem Sohn~\GreStar{}~\nopagebreak

und dem Heiligen Geist.

\noindent Wie im Anfang, so auch jetzt und alle Zeit~\GreStar{}~\nopagebreak

und in Ewigkeit. Amen.

\vspace{10pt}

\noindent \textsc{Ant.} Ich lege mich nieder und ruhe in Frieden. 

\vspace{10pt}

\noindent \textsc{2 Ant.} Mein Leib ruht in sicherer Hoffnung: Du gibst mich der Unterwelt nicht preis.

\subsubsection{Psalm 16, 1--11}

\noindent Behüte mich, Gott, denn ich vertraue dir.~†~\nopagebreak

Ich sage zum Herrn: „Du bist mein Herr;~\GreStar{}~\nopagebreak

mein ganzes Glück bist du allein.“

\noindent An den Heiligen im Lande, den Herrlichen,~\GreStar{}~\nopagebreak

an ihnen nur hab’ ich mein Gefallen.

\noindent Viele Schmerzen leidet, wer fremden Göttern folgt.~†~\nopagebreak

Ich will ihnen nicht opfern,~\GreStar{}~\nopagebreak

ich nehme ihre Namen nicht auf meine Lippen.

\noindent Du, Herr, gibst mir das Erbe und reichst mir den Becher;~\GreStar{}~\nopagebreak

du hältst mein Los in deinen Händen.

\noindent Auf schönem Land fiel mir mein Anteil zu.~\GreStar{}~\nopagebreak

Ja, mein Erbe gefällt mir gut.

\noindent Ich preise den Herrn, der mich beraten hat.~\GreStar{}~\nopagebreak

Auch mahnt mich mein Herz in der Nacht.

\noindent Ich habe den Herrn beständig vor Augen.~\GreStar{}~\nopagebreak

Er steht mir zur Rechten, ich wanke nicht.

\noindent Darum freut sich mein Herz und frohlockt meine Seele;~\GreStar{}~\nopagebreak

auch mein Leib wird wohnen in Sicherheit.

\noindent Denn du gibst mich nicht der Unterwelt preis;~\GreStar{}~\nopagebreak

du lässt deinen Frommen das Grab nicht schauen.

\noindent Du zeigst mir den Pfad zum Leben.~†~\nopagebreak

Vor deinem Angesicht herrscht Freude in Fülle,~\GreStar{}~\nopagebreak

zu deiner Rechten Wonne für alle Zeit.

\noindent Ehre sei dem Vater und dem Sohn~\GreStar{}~\nopagebreak

und dem Heiligen Geist.

\noindent Wie im Anfang, so auch jetzt und alle Zeit~\GreStar{}~\nopagebreak

und in Ewigkeit. Amen.

\vspace{10pt}

\noindent \textsc{Ant.} Mein Leib ruht in sicherer Hoffnung: Du gibst mich der Unterwelt nicht preis.

\vspace{10pt}

\noindent \textsc{3 Ant.} Hebt euch, ihr uralten Pforten! Es kommt der König der Herrlichkeit.

\subsubsection{Psalm 24, 1--10}

\noindent Dem Herrn gehört die Erde und was sie erfüllt,~\GreStar{}~\nopagebreak

der Erdkreis und seine Bewohner.

\noindent Denn er hat ihn auf Meere gegründet,~\GreStar{}~\nopagebreak

ihn über Strömen befestigt.

\noindent Wer darf hinaufziehn zum Berg des Herrn,~\GreStar{}~\nopagebreak

wer darf stehn an seiner heiligen Stätte?

\noindent Der reine Hände hat und ein lauteres Herz,~\GreStar{}~\nopagebreak

der nicht betrügt und keinen Meineid schwört.

\noindent Er wird Segen empfangen vom Herrn~\GreStar{}~\nopagebreak

und Heil von Gott, seinem Helfer.

\noindent Das sind die Menschen, die nach ihm fragen,~\GreStar{}~\nopagebreak

die dein Antlitz suchen, Gott Jakobs.

\noindent Ihr Tore, hebt euch nach oben,~†~\nopagebreak

hebt euch, ihr uralten Pforten;~\GreStar{}~\nopagebreak

denn es kommt der König der Herrlichkeit.

\noindent Wer ist der König der Herrlichkeit?~\GreStar{}~\nopagebreak

Der Herr, stark und gewaltig, der Herr, mächtig im Kampf.

\noindent Ihr Tore, hebt euch nach oben,~†~\nopagebreak

hebt euch, ihr uralten Pforten;~\GreStar{}~\nopagebreak

denn es kommt der König der Herrlichkeit.

\noindent Wer ist der König der Herrlichkeit?~\GreStar{}~\nopagebreak

Der Herr der Heerscharen, er ist der König der Herrlichkeit.

\noindent Ehre sei dem Vater und dem Sohn~\GreStar{}~\nopagebreak

und dem Heiligen Geist.

\noindent Wie im Anfang, so auch jetzt und alle Zeit~\GreStar{}~\nopagebreak

und in Ewigkeit. Amen.

\vspace{10pt}

\noindent \textsc{Ant.} Hebt euch, ihr uralten Pforten! Es kommt der König der Herrlichkeit.

   \subsubsection{Versiculum}  \greannotation{} \index[Versiculum]{Versiculum} \label{Versiculum (Versiculum)} \grecommentary[0pt]{} \gresetinitiallines{0} \grechangestyle{initial}{\fontsize{36}{36}\selectfont} \grechangedim{maxbaroffsettextleft@nobar}{12 cm}{scalable} \grechangedim{spaceabovelines}{0.5cm}{scalable} \gresetlyriccentering{syllable}   \gregorioscore{chants/versiculum_thursday}    \newpage
  \subsection{Lesungen}
   \subsubsection{Erste Lesung} \vspace{5pt} \emph{Aus den Klageliedern des Propheten Jeremia. Aleph. Weh, wie glanzlos ist das Gold, gedunkelt das köstliche Feingold, hingeschüttet die heiligen Steine an den Ecken aller Straßen. Beth. Die edlen Kinder Zions, einst aufgewogen mit reinem Gold, weh, wie Krüge aus Ton sind sie geachtet, wie Werk von Töpferhand. Ghimel. Selbst Schakale reichen die Brust, säugen ihre Jungen. Die Töchter meines Volkes sind grausam wie Strauße in der Wüste. Daleth. Des Säuglings Zunge klebt an seinem Gaumen vor Durst. Kinder betteln um Brot; keiner bricht es ihnen. - Jerusalem, Jerusalem, bekehre dich zum Herrn, deinem Gott.} \vspace{5pt} \greannotation{} \index[Erste Lesung]{Erste Lesung} \label{Erste Lesung (Erste Lesung)} \grecommentary[3pt]{Klgl 4, 1--4} \gresetinitiallines{1} \grechangestyle{initial}{\fontsize{36}{36}\selectfont} \grechangedim{maxbaroffsettextleft@nobar}{12 cm}{scalable} \grechangedim{spaceabovelines}{0.5cm}{scalable} \gresetlyriccentering{vowel}   \gregorioscore{chants/va--lamentationes_71_de_lamentatione--dominican}
   \subsubsection{II} \vspace{5pt} \emph{He. Die einst Leckerbissen schmausten, verschmachten auf den Straßen. Die einst auf Purpur lagen, wälzen sich jetzt im Unrat. Vau. Größer ist die Schuld der Tochter, meines Volkes, als die Sünde Sodoms, das plötzlich vernichtet wurde, ohne dass eine Hand sich rührte. Zain. Ihre jungen Männer waren reiner als Schnee, weißer als Milch, ihr Leib rosiger als Korallen, saphirblau ihre Adern. Heth. Schwärzer als Ruß sehen sie aus, man erkennt sie nicht auf den Straßen. Die Haut schrumpft ihnen am Leib, trocken wie Holz ist sie geworden. - Jerusalem, Jerusalem, bekehre dich zum Herrn, deinem Gott.} \vspace{5pt} \greannotation{} \index[II]{II} \label{II (II)} \grecommentary[3pt]{Klgl 4, 5--8} \gresetinitiallines{1} \grechangestyle{initial}{\fontsize{36}{36}\selectfont} \grechangedim{maxbaroffsettextleft@nobar}{12 cm}{scalable} \grechangedim{spaceabovelines}{0.5cm}{scalable} \gresetlyriccentering{vowel}   \gregorioscore{chants/va--lamentationes_72_he_qui_vescebantur--dominican}
   \subsubsection{III} \vspace{5pt} \emph{Teth. Besser die vom Schwert Getöteten als die vom Hunger Getöteten; sie sind verschmachtet, vom Missertrag der Felder getroffen. Jod. Die Hände liebender Mütter kochten die eigenen Kinder. Sie dienten ihnen als Speise beim Zusammenbruch der Tochter, meines Volkes. Caph. Randvoll gemacht hat der Herr seinen Grimm, ausgegossen seinen glühenden Zorn. Er entfachte in Zion ein Feuer, das bis auf den Grund alles verzehrte. Lamed. Kein König eines Landes, kein Mensch auf der Erde hätte jemals geglaubt, dass ein Bedränger und Feind durchschritte die Tore Jerusalems. - Jerusalem, Jerusalem, bekehre dich zum Herrn, deinem Gott.}  \vspace{5pt} \greannotation{} \index[III]{III} \label{III (III)} \grecommentary[3pt]{Klgl 4, 9--12} \gresetinitiallines{1} \grechangestyle{initial}{\fontsize{36}{36}\selectfont} \grechangedim{maxbaroffsettextleft@nobar}{12 cm}{scalable} \grechangedim{spaceabovelines}{0.5cm}{scalable} \gresetlyriccentering{vowel}   \gregorioscore{chants/va--lamentationes_73_teth_melius_fuit--dominican}    \newpage
   \subsubsection{Zweite Lesung}             \vspace{10pt}

Leo der Große († 461), Aus einer Predigt über die Passion des Herrn

\vspace{10pt}

%\par \hfill(Nn. 65071: SC 123, 95--101)

\lettrine[lines=3]{M}{}it allen Kräften unseres Geistes und unseres Leibes müssen wir darnach trachten, unzertrennlich mit dem Geheimnis des Leidens Christi verbunden zu bleiben; denn der Herr sagt: „Wer nicht sein Kreuz auf sich nimmt und mir nachfolgt, ist meiner nicht würdig.“ Und der Apostel spricht: „Wenn wir mit ihm leiden, werden wir auch mit ihm verherrlicht werden.“ Wer anders erweist also nach diesen Worten dem wahrhaft leidenden, sterbenden und auferstehenden Christus seine Verehrung, als wer mit ihm leidet, stirbt und aufersteht? Diese Teilnahme an dem Leiden des Herrn hat bei allen Kindern der Kirche schon mit ihrer wunderbaren Wiedergeburt begonnen: Durch die Tilgung der Sünde ersteht hier der Mensch zu neuem Leben, und durch das dreimalige Untertauchen wird der dreitägige Tod des Herrn versinnbildet. Bei der Taufe wird gleichsam die Erddecke von einem Grabe entfernt. Mit unserem alten Menschen steigen wir in den Taufquell hinab, und neugeboren kommen wir aus ihm hervor. Was aber durch dieses Sakrament mit uns begonnen wurde, das müssen wir durch Taten vollenden. Die ganze Lebenszeit, die den im Heiligen Geiste Wiedergeborenen noch übrigbleibt, muß ein beständiges Tragen des Kreuzes sein. Obgleich nämlich durch die Macht des Leidens Christi dem starken und grausamen Feinde unseres Geschlechtes die „Gefäße der alten Erbeutung“ entrissen wurden und „der Herrscher dieser Welt“ über die Herzen der Erlösten keine Gewalt mehr hat, verfolgt er doch die Menschen selbst nach ihrer Rechtfertigung immer noch mit seiner alten Bosheit. Auf mancherlei Art greift er die an, in denen er nicht mehr herrscht, um nachlässige und sorglose Seelen aufs neue mit noch grausameren Banden an sich zu ketten, um sie aus dem Paradies der Kirche zu vertreiben und sie zu Genossen seiner Verdammnis zu machen. Wenn darum jemand merkt, daß er die Grenzen der christlichen Gebote überschreitet und daß seine Begierden auf Dinge gerichtet sind, die ihn vom rechten Wege abbringen könnten, so nehme er seine Zuflucht zum Kreuze des Herrn und kreuzige sein sündhaftes Wollen und Wünschen auf dem Baume des Lebens.
    \newpage
   \subsubsection{Oratio Ieremiae Prophetae} \vspace{5pt} \emph{Gebet des Propheten Jeremia. Herr, denk daran, was uns geschehen, blick her und sieh unsre Schmach! An Ausländer fiel unser Erbe, unsre Häuser kamen an Fremde. Wir wurden Waisen, Kinder ohne Vater, unsere Mütter wurden Witwen. Unser Wasser trinken wir für Geld, unser Holz müssen wir bezahlen. Wir werden getrieben, das Joch auf dem Nacken, wir sind müde, man versagt uns die Ruhe. Nach Ägypten streckten wir die Hand, nach Assur, um uns mit Brot zu sättigen. Unsere Väter haben gesündigt; sie sind nicht mehr. Wir müssen ihre Sünden tragen. Sklaven herrschen über uns, niemand entreißt uns ihren Händen. Unter Lebensgefahr holen wir unser Brot, bedroht vom Schwert der Wüste. Unsere Haut glüht wie ein Ofen von den Gluten des Hungers. Frauen hat man in Zion geschändet, Jungfrauen in den Städten von Juda. Fürsten wurden von Feindeshand gehängt, den Ältesten nahm man die Ehre. Junge Männer mussten die Handmühlen schleppen, unter der Holzlast brachen Knaben zusammen. Die Alten blieben fern vom Tor, die Jungen vom Saitenspiel. Dahin ist unseres Herzens Freude, in Trauer gewandelt unser Reigen. Die Krone ist uns vom Haupt gefallen. Weh uns, wir haben gesündigt. Darum ist krank unser Herz, darum sind trüb unsere Augen über den Zionsberg, der verwüstet liegt; Füchse laufen dort umher. Du aber, Herr, bleibst ewig, dein Thron von Geschlecht zu Geschlecht. Warum willst du uns für immer vergessen, uns verlassen fürs ganze Leben? Kehre uns, Herr, dir zu, dann können wir uns zu dir bekehren. Erneuere unsere Tage, damit sie werden wie früher. Oder hast du uns denn ganz verworfen, zürnst du uns über alle Maßen? - Jerusalem, Jerusalem, bekehre dich zum Herrn, deinem Gott.}  \vspace{5pt} \greannotation{} \index[Oratio Ieremiae Prophetae]{Oratio Ieremiae Prophetae} \label{Oratio Ieremiae Prophetae (Oratio Ieremiae Prophetae)} \grecommentary[4pt]{Klgl 5, 1--22} \gresetinitiallines{1} \grechangestyle{initial}{\fontsize{36}{36}\selectfont} \grechangedim{maxbaroffsettextleft@nobar}{12 cm}{scalable} \grechangedim{spaceabovelines}{0.5cm}{scalable} \gresetlyriccentering{vowel}   \gregorioscore{chants/misc.oratio.ieremiae}    \newpage
  \section{Laudes}
  \subsubsection{Psalmodie} \noindent \textsc{1 Ant.} Sie klagen um ihn, wie man klagt um den einzigen Sohn; denn er wurde getötet - und war doch ohne Schuld.

\subsubsection{Psalm 64, 2-11}

\noindent Höre, o Gott, mein lautes Klagen,~\GreStar{}~\nopagebreak

schütze mein Leben vor dem Schrecken des Feindes!

\noindent Verbirg mich vor der Schar der Bösen,~\GreStar{}~\nopagebreak

vor dem Toben derer, die Unrecht tun.

\noindent Sie schärfen ihre Zunge wie ein Schwert,~\GreStar{}~\nopagebreak

schießen giftige Worte wie Pfeile,

\noindent um den Schuldlosen von ihrem Versteck aus zu treffen.~\GreStar{}~\nopagebreak

Sie schießen auf ihn, plötzlich und ohne Scheu.

\noindent Sie sind fest entschlossen zu bösem Tun.~\GreStar{}~\nopagebreak

Sie planen, Fallen zu stellen, und sagen: „Wer sieht uns schon?“

\noindent Sie haben Bosheit im Sinn,~\GreStar{}~\nopagebreak

doch halten sie ihre Pläne geheim. 

\noindent Ihr Inneres ist heillos verdorben,~\GreStar{}~\nopagebreak

ihr Herz ist ein Abgrund.

\noindent Da trifft sie Gott mit seinem Pfeil;~\GreStar{}~\nopagebreak

sie werden jählings verwundet.

\noindent Ihre eigene Zunge bringt sie zu Fall.~\GreStar{}~\nopagebreak

Alle, die es sehen, schütteln den Kopf.

\noindent Dann fürchten sich alle Menschen;~†~\nopagebreak

sie verkünden Gottes Taten~\GreStar{}~\nopagebreak

und bedenken sein Wirken.

\noindent Der Gerechte freut sich am Herrn und sucht bei ihm Zuflucht.~\GreStar{}~\nopagebreak

Und es rühmen sich alle Menschen mit redlichem Herzen.

\noindent Ehre sei dem Vater und dem Sohn~\GreStar{}~\nopagebreak

und dem Heiligen Geist.

\noindent Wie im Anfang, so auch jetzt und alle Zeit~\GreStar{}~\nopagebreak

und in Ewigkeit. Amen.

\vspace{10pt}

\noindent \textsc{Ant.} Sie klagen um ihn, wie man klagt um den einzigen Sohn; denn er wurde getötet - und war doch ohne Schuld.

%\vspace{10pt} 

\newpage

\noindent \textsc{2 Ant.} Vor den Pforten der Unterwelt rette, o Herr, mein Leben.

\subsubsection{Jes 38, 10-13a.14c-d.17-20}

\noindent Ich sagte: In der Mitte meiner Tage~†~\nopagebreak

muss ich hinab zu den Pforten der Unterwelt,~\GreStar{}~\nopagebreak

man raubt mir den Rest meiner Jahre.

\noindent Ich darf den Herrn nicht mehr schauen im Land der Lebenden,~\GreStar{}~\nopagebreak

keinen Menschen mehr sehen bei den Bewohnern der Erde.

\noindent Meine Hütte bricht man über mir ab,~\GreStar{}~\nopagebreak

man schafft sie weg wie das Zelt eines Hirten. 

\noindent Wie ein Weber hast du mein Leben zu Ende gewoben,~\GreStar{}~\nopagebreak

du schneidest mich ab wie ein fertig gewobenes Tuch. 

\noindent Vom Anbruch des Tages bis in die Nacht gibst du mich preis;~\GreStar{}~\nopagebreak

bis zum Morgen schreie ich um Hilfe.

\noindent Meine Augen blicken ermattet nach oben:~\GreStar{}~\nopagebreak

Ich bin in Not, Herr. Steh mir bei!

\noindent Du hast mich aus meiner bitteren Not gerettet,~†~\nopagebreak

du hast mich vor dem tödlichen Abgrund bewahrt;~\GreStar{}~\nopagebreak

denn all meine Sünden warfst du hinter deinen Rücken.

\noindent Ja, in der Unterwelt dankt man dir nicht,~†~\nopagebreak

die Toten loben dich nicht;~\GreStar{}~\nopagebreak

wer ins Grab gesunken ist, kann nichts mehr von deiner Güte erhoffen.

\noindent Nur die Lebenden danken dir, wie ich am heutigen Tag.~\GreStar{}~\nopagebreak

Von deiner Treue erzählt der Vater den Kindern.

\noindent Der Herr war bereit, mir zu helfen.~\GreStar{}~\nopagebreak

Wir wollen singen und spielen im Haus des Herrn, solange wir leben!

\noindent Ehre sei dem Vater und dem Sohn~\GreStar{}~\nopagebreak

und dem Heiligen Geist.

\noindent Wie im Anfang, so auch jetzt und alle Zeit~\GreStar{}~\nopagebreak

und in Ewigkeit. Amen.

\vspace{10pt}

\noindent \textsc{Ant.} Vor den Pforten der Unterwelt rette, o Herr, mein Leben.

%\vspace{10pt} 

\newpage

\noindent \textsc{3 Ant.} Ich war tot, doch ich lebe in Ewigkeit. Ich habe die Schlüssel des Todes und der Unterwelt.

\subsubsection{Psalm 150, 1-6}

\noindent Lobet Gott in seinem Heiligtum,~\GreStar{}~\nopagebreak

lobt ihn in seiner mächtigen Feste!

\noindent Lobt ihn für seine großen Taten,~\GreStar{}~\nopagebreak

lobt ihn in seiner gewaltigen Größe!

\noindent Lobt ihn mit dem Schall der Hörner,~\GreStar{}~\nopagebreak

lobt ihn mit Harfe und Zither!

\noindent Lobt ihn mit Pauken und Tanz,~\GreStar{}~\nopagebreak

lobt ihn mit Flöten und Saitenspiel!

\noindent Lobt ihn mit hellen Zimbeln,~\GreStar{}~\nopagebreak

lobt ihn mit klingenden Zimbeln!

\noindent Alles, was atmet,~\GreStar{}~\nopagebreak

lobe den Herrn!

\noindent Ehre sei dem Vater und dem Sohn~\GreStar{}~\nopagebreak

und dem Heiligen Geist.

\noindent Wie im Anfang, so auch jetzt und alle Zeit~\GreStar{}~\nopagebreak

und in Ewigkeit. Amen.

\vspace{10pt}

\noindent \textsc{Ant.} Ich war tot, doch ich lebe in Ewigkeit. Ich habe die Schlüssel des Todes und der Unterwelt.
                 \newpage
   \subsubsection{Kurzlesung}     \hfill Hos 6, 1--2        \lettrine[lines=3]{K}{}ommt, wir kehren zum Herrn zurück! Denn er hat Wunden gerissen, er wird uns auch heilen; er hat verwundet, er wird auch verbinden. Nach zwei Tagen gibt er uns das Leben zurück, am dritten Tag richtet er uns wieder auf, und wir leben vor seinem Angesicht.
   \subsubsection{Responsorium} \emph{Christus war für uns gehorsam bis zum Tod, bis zum Tod am Kreuze. \Vbar.~Darum auch hat Gott ihn erhöht und ihm den Namen gegeben, der über allen Namen steht.} \vspace{5pt}  \greannotation{V} \index[Responsorium]{Responsorium} \label{Responsorium (Responsorium)} \grecommentary[7pt]{Phil 2, 8; \Vbar.~9} \gresetinitiallines{1} \grechangestyle{initial}{\fontsize{36}{36}\selectfont} \grechangedim{maxbaroffsettextleft@nobar}{12 cm}{scalable} \grechangedim{spaceabovelines}{0.5cm}{scalable} \gresetlyriccentering{vowel}   \gregorioscore{chants/gr--christus_factus_est--dominican--id_6737}    \newpage
  \subsubsection{Benedictus} \noindent \textsc{Ant.} Retter der Welt, errette uns! Du hast uns erlöst durch dein Kreuz und dein Blut. Hilf uns, Herr, unser Gott!

\paragraph{Lk 1, 68-79}

\gabcsnippet{(f3) (e f hr0 fR ,3) (hr0 iRr1 fR : hr0 e f gRr1 fR ::)}

\noindent Gepriesen sei der Herr, der Gott \uline{I}sraels!~\GreStar{}~\nopagebreak

Denn er hat sein Volk besucht und ihm Erlös\uline{u}ng geschaffen;

\noindent er hat uns einen starken Retter erw\uline{e}ckt~\GreStar{}~\nopagebreak

im Hause seines Kn\uline{e}chtes David.

\noindent So hat er verheißen von \uline{a}lters her~\GreStar{}~\nopagebreak

durch den Mund seiner heilig\uline{e}n Propheten.

\noindent Er hat uns errettet vor unseren F\uline{ei}nden~\GreStar{}~\nopagebreak

und aus der Hand aller, d\uline{ie} uns hassen;

\noindent er hat das Erbarmen mit den Vätern an uns vollend\uline{e}t~†~\nopagebreak

und an seinen heiligen B\uline{u}nd gedacht~\GreStar{}~\nopagebreak

an den Eid, den er unserm Vater Abrah\uline{a}m geschworen hat;

\noindent er hat uns geschenkt, dass wir, aus Feindeshand befr\uline{ei}t,~†~\nopagebreak

ihm furchtlos dienen in Heiligkeit und Ger\uline{e}chtigkeit~\GreStar{}~\nopagebreak

vor seinem Angesicht all \uline{u}nsre Tage.

\noindent Und du, Kind, wirst Prophet des Höchst\uline{en hei}ßen;~†~\nopagebreak

denn du wirst dem Herrn vor\uline{a}ngehn~\GreStar{}~\nopagebreak

und ihm den W\uline{e}g bereiten.

\noindent Du wirst sein Volk mit der Erfahrung des Heils besch\uline{e}nken~\GreStar{}~\nopagebreak

in der Vergeb\uline{u}ng der Sünden. 

\noindent Durch die barmherzige Liebe unseres G\uline{o}ttes~\GreStar{}~\nopagebreak

wird uns besuchen das aufstrahlende Licht \uline{au}s der Höhe,

\noindent um allen zu leuchten, die in Finsternis sitzen und im Schatten des T\uline{o}des,~\GreStar{}~\nopagebreak

und unsre Schritte zu lenken auf den W\uline{e}g des Friedens.

\noindent Ehre sei dem Vater und dem S\uline{o}hn~\GreStar{}~\nopagebreak

und dem H\uline{ei}ligen Geist.

\noindent Wie im Anfang, so auch jetzt und \uline{a}lle Zeit~\GreStar{}~\nopagebreak

und in Ew\uline{i}gkeit. Amen.

\vspace{5pt}

\noindent \textsc{Ant.} Retter der Welt, errette uns! Du hast uns erlöst durch dein Kreuz und dein Blut. Hilf uns, Herr, unser Gott!
                 \newpage
   \subsubsection{Preces}  \greannotation{} \index[Preces]{Holy Thursday} \label{Holy Thursday (Preces)} \grecommentary[3pt]{} \gresetinitiallines{1} \grechangestyle{initial}{\fontsize{36}{36}\selectfont} \grechangedim{maxbaroffsettextleft@nobar}{12 cm}{scalable} \grechangedim{spaceabovelines}{0.7cm}{scalable} \gresetlyriccentering{vowel}   \gregorioscore{chants/misc.versus_litanici_in_cantu_sabbato_sancto-de-rubrics}
   \subsubsection{Vater Unser}   \index[Vater Unser]{Vater Unser} \label{Vater Unser (Vater Unser)}             \vspace{10pt}
   \subsubsection{Oration}   \index[Oration]{Oration} \label{Oration (Oration)}         \lettrine[lines=3]{A}{}llmächtiger, ewiger Gott, dein eingeborener Sohn ist in das Reich des Todes hinabgestiegen und von den Toten glorreich auferstanden. Gib, dass deine Gläubigen, die durch die Taufe mit ihm begraben wurden, durch seine Auferstehung zum ewigen Leben gelangen. Darum bitten wir durch ihn, Jesus Christus, deinen Sohn, unseren Herrn und Gott, der in der Einheit des Heiligen Geistes mit dir lebt und herrscht in alle Ewigkeit.
\par \Rbar.~Amen.

   \subsubsection{Schlußsegen}   \index[Schlußsegen]{Schlußsegen} \label{Schlußsegen (Schlußsegen)}         Der Herr sei mit euch.

\Rbar. Und mit deinem Geiste.

Es segne euch der allmächtige Gott, + \par der Vater und der Sohn und der Heilige Geist.

\Rbar. Amen.

Gehet hin in Frieden.

\Rbar. Dank sei Gott, dem Herrn.


  \end{document}
